\documentclass[12pt]{article}
\title{Digest of The New York Times}
\author{The New York Times}

% \usepackage[pdftex,twocolumn]{geometry}
% \geometry{paperwidth=34.15cm}
% \geometry{paperheight=23.04cm}
% \geometry{margin=11mm}
% \geometry{columnsep=8mm}


% \definecolor{dark-green}{rgb}{0,0.25,0}
% \usepackage[utf8]{inputenc}

\usepackage{fontspec,xunicode,xltxtra}
\setmainfont[Ligatures=Required,Mapping=tex-text]{Georgia}
% \setmainfont[Ligatures=Required,Mapping=tex-text]{Palatino Linotype}
\setsansfont[Ligatures=Required,Mapping=tex-text]{Myriad Pro}
\setmonofont[Ligatures=Required,Mapping=tex-text]{Courier New}
\usepackage{xeCJK}
\setCJKmainfont{宋体}
\setCJKmonofont{黑体}

% \usepackage{indentfirst}


\usepackage[dvips,dvipsnames,svgnames]{xcolor}
\definecolor{light-gray}{gray}{0.95}

%生成PDF的链接
\usepackage[dvipdfmx]{hyperref}
\hypersetup{
    bookmarks=true,         % show bookmarks bar?
    bookmarksopen=true,
    pdfpagemode=UseNone,    % options: UseNode, UseThumbs, UseOutlines, FullScreen
    pdfstartview=FitH,
    pdfborder=1,
    pdfhighlight=/P,
    pdfauthor={wuxch},
    unicode=true,           % non-Latin characters in Acrobat’s bookmarks
    colorlinks,             % false: boxed links; true: colored links
    linkcolor=blue,         % color of internal links
    citecolor=blue,        % color of links to bibliography
    filecolor=magenta,      % color of file links
    urlcolor=cyan           % color of external links
}
% 使用目录
\usepackage{makeidx}
% 使用图形
\usepackage{graphicx}


\usepackage{titletoc}
\dottedcontents{section}[1.5em]{}{2.3em}{1pc}

% % 页眉页脚的设定
% \usepackage{fancyhdr}
% \pagestyle{fancy}
% \lhead{\includegraphics[height=0.2in]{zte_logo}}

% \usepackage{footmisc}
\usepackage{pfnote}
\usepackage{fnpos}

% % 嵌入的代码显示
% \usepackage{listings}
% \lstset{language=C++, breaklines, extendedchars=false}
% \lstset{keywordstyle=\color{blue},
%         commentstyle=\color{SeaGreen},
%         stringstyle=\ttfamily,
%         showstringspaces=false,
%         backgroundcolor=\color{light-gray}}
% \lstset{basicstyle=\small}

\usepackage{booktabs}

\usepackage{tipa}
\usepackage{phonetic}
\usepackage[gen]{eurosym}

% \usepackage{CJK}
\usepackage{makeidx}
% \usepackage{glossary}

\usepackage{fancyhdr}

\usepackage{lastpage}
\cfoot{\textsf{\thepage\ of \pageref*{LastPage}}}

\pagestyle{fancy}
\fancyhead{}
\renewcommand{\headrulewidth}{0pt}
\pagenumbering{arabic}

\usepackage{graphicx}
\usepackage{picinpar}

\setlength{\parindent}{0in}
\setlength{\parskip}{2.5ex plus4mm minus3mm}
% \setlength{\parskip}{.3\baselineskip}

\setcounter{section}{0}

% bibtex
\usepackage[super]{natbib}
% \usepackage[globalcitecopy,labelstoglobalaux]{bibunits}
% \makeatletter
% \renewcommand{\@bibunitname}{\jobname.\the\@bibunitauxcnt}
% \makeatother

% \renewcommand{\refname}{缩略语}
% \renewcommand{\refname}{}

\makeatletter                   % Reference list option change
\renewcommand\@biblabel[1]{#1:  } % from [1] to 1: 
\makeatother                    %

\newcommand{\citenumfont}[1]{\textsf{#1}}

% \usepackage{tipa}
% \usepackage{phonetic}

\newcommand{\phonetic}[1]{{\fontspec[Colour=CC00CC]{Microsoft Sans Serif} #1}}


%%% Local Variables:
%%% TeX-master: "medical_trivia"
%%% End:


\makeindex
\begin{document}
\date{}
\thispagestyle{empty}
\renewcommand\contentsname{\textsf{Digest of The New York Times}}
{\footnotesize\textsf{\tableofcontents}}
\clearpage
\setcounter{page}{1}
\section{Ideas Online, Yes, but Some Not So Presidential}

WASHINGTON -- On Jan.~21, his first full day in office, President Obama promised to open up the
government, ordering officials to use modern technologies like Internet message boards and blogs to
give all Americans a bigger voice in public policy.

Well, the people have spoken. But many of them are not sticking to the topics at hand.

The White House made its first major entree into government by the people last month when it set up
an online forum to ask ordinary people for their ideas on how to carry out the president's
open-government pledge. It got an earful -- on legalizing marijuana, revealing U.F.O. secrets and
verifying Mr.~Obama's birth certificate to prove he was really born in the United States and thus
eligible to be president.

``Please, as fellow human beings of this great planet Earth, disclose all known information on
space/UFO's because the world needs to know,'' wrote sprinter5160 on the site, whitehouse.gov/open,
which attracted thousands of similar comments on fringe topics.

While it was not exactly what administration officials had in mind, they noted that democracy can be
a bit messy.

``Even for people who want to talk about U.F.O.'s or the Kennedy assassination, we have created a
forum for people to have a conversation with each other, and potentially to go off and organize and
develop this further,'' said Beth Simone Noveck, a New York Law School professor who is Mr.~Obama's
deputy chief technology officer for open government.

Most of the suggestions were closely related to the topic at hand, like publishing a list of
everyone who meets with the president, using computer graphics to track how rapidly agencies respond
to Freedom of Information Act requests and installing webcams to monitor federal offices. The
administration's goal is to devise regulations that would tell federal agencies how to make their
operations more open to the public.

Undeterred by some of the wilder suggestions, the White House proceeded Monday with the third phase
of the process -- asking people to collaborate online to draft language that could be used to create
the final rules.

The experience so far shows just how hard it is to allow all voices to be heard and still have a
coherent discussion. When millions of Internet users are invited to discuss every regulation, how
can any real work get done? On the other hand, why bother opening up the government if views that
are outside the mainstream -- as defined by the usual collection of lobbyists and think tank
scholars -- are summarily dismissed?

The responsibility for sorting it all out falls to Ms.~Noveck. She has permitted any proposal that
was not abusive or repetitive onto the brainstorming site, just as the Obama transition team did not
stop visitors to its Change.Gov site last fall from voting marijuana legalization as their top
concern for the president-elect.

She argues that the experience of collaborative Web sites like Wikipedia proves that groups of users
can police sites to keep small groups from spoiling things for everyone else. During the public
brainstorming about rules for open government, the White House asked visitors to vote on the best
ideas by clicking a thumbs-up or thumbs-down button, much as people vote on the most interesting
news articles on sites like Digg.

The visitors advanced more than 3,900 ideas, which in turn spawned 11,000 comments that received
210,000 thumb votes.

The result? Three of the top 10 most popular ideas called for legalizing marijuana, and two featured
conspiracy theories about Mr.~Obama's true place of birth. (Among the technical ideas that got a lot
of support was a proposal to have the federal government press states and cities to follow
open-government principles and a call for a central Internet site for all requests to the president
and Congress, modeled after a site for petitions to the British prime minister.)

``This is Obama's Madisonian moment,'' said Clay Shirky, a professor at New York University and the
author of ``Here Comes Everybody,'' a book about Internet collaboration. Just as James Madison, the
nation's fourth president, argued during the drafting of the Constitution that the government must
protect the minority against the tyranny of the majority, Mr.~Shirky said that government must also
prevent small groups of loudmouths from hijacking the public debate.

``The first thing that happens when my mom and dad log into the system and they find it's populated
by U.F.O. people and birth-certificate people, they simply are not going to participate,'' he said.

The White House tried to screen out some of the more unusual comments in the second phase of the
process. Ms.~Noveck summarized the most significant ideas, then invited comments on them at
blog.ostp.gov. Visitors could flag off-topic comments, which were then shunted to a separate part of
the site. That reduced the birth-certificate and U.F.O. comments to a relative trickle.

On Monday, the White House began Phase 3 of its project using yet another format: a wiki, an online
tool that allows a group of people to collectively create and edit documents. Visitors will be able
to submit and edit drafts of the open-government rules, similar to how people contribute to
Wikipedia, the user-created online encyclopedia.

Ultimately, of course, ``this is not policy by referendum,'' said Ms.~Noveck. The Office of
Management and Budget will consider the public comments and the views of agency officials and White
House staff and then put together its own formal draft of the open-government rules. After
soliciting another round of public comment, the final rules will become effective and will govern
the actions of federal agencies.

To some, the bumps in the process simply represent growing pains for a new and promising approach to
government.

``The U.F.O. thing is healthy,'' said Micah L.~Sifry, the editor of TechPresident.com, a blog on
politics and the Internet. ``It's weird there are some groups of people obsessed with it, but it's a
democracy, and you can't make them go away.''

As people get used to this kind of participation, he said, ``the mischief will be much less
noticeable.''

Ms.~Noveck has some confidence that the effort will result in better government because she has built something like this before. As a professor, she worked with the United States Patent Office to test a system that invited the public to help evaluate patent applications. Companies that apply for a lot of patents, like I.B.M. and General Electric, participated in the optional program because the public comments helped patent examiners consider their applications more quickly.

But then, I.B.M. never tried to patent a U.F.O.

\section{Sanford Case a New Dose of Bad News for Republicans}

WASHINGTON -- Republicans were just starting to breathe a little easier.

The news that Senator John Ensign had had an affair with a former aide who was married to another
former aide was fading. Polls showed some voter impatience with President Obama's policies, if not
with the president himself. And the Politico, the insidery Web site that is widely read in the
capital's political precincts, even featured an article exploring the possibility of a Republican
Party comeback.

Then Gov. Mark Sanford of South Carolina, a fiscal conservative seen by many Republicans as an
attractive standard-bearer for the next presidential campaign, went missing. Worse, he returned.

His confession on Wednesday that he had been in Argentina with a woman not his wife -- and not
hiking the Appalachian Trail as his staff had said Monday -- was another jolt of bad news for a
party that has struggled to get off the ropes all year.

That it was the second such confession in little more than a week from a potential Republican
presidential contender -- Mr.~Ensign had been exploring a run in 2012 as well -- left party leaders
dazed. They spent Wednesday alternating between gallows humor and yet another round of conversations
about what the party stands for and who will give it its best shot to retake the White House.

``Personal circumstances over the course of the last week have managed to shrink the front line of
the 2012 possible-contender list by 30 percent,'' said Phil Musser, a former executive director of
the Republican Governors Association.

Speaking of Mr.~Sanford's confession, Mr.~Musser said, ``The concern here is that this continues a
broader narrative that is completely unhelpful to the Republican Party's rebuilding -- that's life,
but it's a personal tragedy that fairly or unfairly compounds a series of problems.''

That series of problems has become so chronic that even the party's most pragmatic members could be
forgiven for wondering whether being named ``possible 2012 contender'' is like winning the movie
role of Superman, long believed by some to carry a curse for those actors who don his blue tights.

One by one, those who have been publicly discussed as possible Republican candidates in 2012 have
stumbled.

Gov. Bobby Jindal of Louisiana suffered a political setback after even his fellow conservatives
harshly critiqued his televised response to Mr.~Obama's prime-time address to Congress in February.
The speech, which was supposed to provide a moment to shine in front of a national audience, instead
became fodder for late-night comedy.

Gov. Sarah Palin of Alaska, the former Republican vice-presidential nominee who was eviscerated by
some of her own political aides at the end of last year's presidential race, continued to get
national attention, but hardly the kind likely to help convince voters that she would be a
substantive candidate. The father of her unwed teenage daughter's baby feuded openly with the Palin
family, and the governor exasperated some Republicans in Washington with her off-again, on-again
plans for headlining a fund-raiser there.

After basking in glowing reviews among political pundits this year, Newt Gingrich, the former House
speaker, had to apologize for a post on Twitter in which he called Mr.~Obama's Supreme Court
nominee, Sonia Sotomayor, ``racist'' for saying that she hoped Latinas would be generally better
equipped to make judicial decisions than their white male counterparts.

Another possible Republican presidential candidate in 2012, Gov. Jon M.~Huntsman Jr.~of Utah, fell
out of contention when he accepted Mr.~Obama's offer to become ambassador to China, robbing the
party of a rising star.

All of their troubles have served to improve the prospects of other contenders who have generally
stayed out of the spotlight this year, or have ventured into it only gingerly, like former Gov. Mitt
Romney of Massachusetts, Gov. Tim Pawlenty of Minnesota and Gov. Haley Barbour of Mississippi.

Some prominent party members argued that criticism in the mainstream news media of Ms.~Palin,
Mr.~Gingrich and Mr.~Jindal did not reflect their standing among the conservative voters who decide
primaries and caucuses -- and that the confessions of Mr.~Ensign and Mr.~Sanford would be viewed in
isolation.

``I disagree with the idea that this shows problems for the modern Republican Party,'' said Grover
Norquist, the president of Americans for Tax Reform, a group that applauded Mr.~Sanford's attempt to
refuse some federal stimulus funds earlier this year. In reference to the fiscally conservative
philosophies of Mr.~Ensign and Mr.~Sanford, he joked, ``I think instead it shows that sexual
attractiveness of limited-government conservatism.''

As television pundits noted on Wednesday, confessions by former Gov. Eliot Spitzer of New York that
he had been involved with a prostitute and by former Gov. Jim McGreevey of New Jersey that he had
been unfaithful to his wife with a gay lover did not hurt Democrats nationally, although both men
resigned.

But other senior Republican strategists and leaders said they were concerned that their party's
large segment of evangelical voters makes the party more vulnerable to political damage from
scandal, especially when it involves politicians like Mr.~Sanford and Mr.~Ensign, who had both been
harshly critical of the infidelities of former President Bill Clinton and others.

``When we do these kinds of things like what happened with Ensign and now with Sanford it hurts our
credibility as a party of good governing and of values,'' said Ron Kaufman, a Republican lobbyist
who is close to Mr.~Romney. Mr.~Kaufman is among those in his party who believe that the news that
former Representative Mark Foley of Florida had sent sexually explicit e-mail messages to male
Congressional pages cost the party in 2006 and 2008.

``I think there is somewhat of an identity crisis in the Republican Party,'' said Tony Perkins,
president of the Family Research Council, an evangelical group in Washington. ``Are they going to be
a party that attracts values voters, and are they going to be the party that lives by those
values?''

\section{Obama Set for First Step on Immigration Reform}

WASHINGTON -- President Obama is expected to meet with Congressional leaders of both parties on
Thursday to begin laying the political groundwork for sweeping immigration legislation, even though
its passage this year is considered very unlikely.

With lawmakers already immersed in health care, financial regulation and energy policy, and with the
Senate set to hold hearings soon on Judge Sonia Sotomayor's nomination to the Supreme Court,
administration officials and many in Congress say it is improbable that they will be able to add
anything as challenging as an immigration overhaul.

And the clock is not the only obstacle. While there is a consensus that the immigration system is
broken, Republicans and Democrats, politically burned over the issue in the recent past, remain
divided even within their own parties over how to fix it.

The unemployment rate is expected to stay high, making Democrats who are wavering on immigration
reform leery of supporting it. And while polls show that Hispanic voters care deeply about changing
the system -- many of them are related to or at least know someone who is living in this country
illegally -- even they see it as a lower priority than the economy and health care.

Then there is the question of whether it is Congress or the White House that will take the lead.

Senator Edward M.~Kennedy, Democrat of Massachusetts and a longtime proponent of expanded guest
worker programs and legal status for the estimated 12 million illegal immigrants currently in the
country, has been succeeded as chairman of the Senate Judiciary subcommittee on immigration by
Senator Charles E.~Schumer, Democrat of New York. In that role, Mr.~Schumer would take the point in
pushing for passage of a new bill.

But Republicans have refused to put their political capital at risk without some assurances that
Mr.~Obama will spend some of his own. Senator John Cornyn of Texas, the senior Republican on
Mr.~Schumer's subcommittee, said: ``So far what we have seen from the White House, frankly, is a lot
of photo-ops and not a lot of rolling up your sleeves and doing the hard work. And this is an issue
that is going to take a lot of hard work.''

Aides to Mr.~Obama say he does not intend to get out in front of any proposal until there is a
strong bipartisan commitment to pass it. That stance has the potential to paralyze the process,
since lawmakers are looking to him to use his bully pulpit, and high approval ratings, to help them
fend off any political backlash among their constituents.

``His position is very clear: he thinks we need comprehensive immigration reform,'' David Axelrod, a
senior adviser to the president, said in an interview. ``But that's not something that's going to
happen simply on his volition.''

``Obviously work needs to be done,'' Mr.~Axelrod added, ``and not just from our end, but from the
proponents in Congress, to bring it to the point where it can get passed.''

How to move Congress to that point is likely to be the focus of Thursday's meeting at the White
House, lawmakers and presidential aides say.

In recent days, the Senate majority leader, Harry Reid, Democrat of Nevada, has said he believes
there are enough votes to pass a comprehensive bill this year. And in a speech here Wednesday before
the Migration Policy Institute, Mr.~Schumer agreed, though he said the key to assuaging opponents
was to show that any new immigration legislation would not only legalize the status of illegal
immigrants already in the country but also include tough measures to prevent new waves from
entering.

Public sentiment about the proposed immigration overhaul that failed in the Senate two years ago was
that it ``was too soft on illegal immigrants,'' Mr.~Schumer said, adding, ``Unless we can convince
Americans we're going to be really tough, then this is not going to work.''

If the outlook is unclear in the Senate, it is even more so in the House, where at least 40
Democrats represent conservative districts. Among them is Representative Heath Shuler of North
Carolina, who says he supports enforcing current immigration laws at the border and in the workplace
before considering comprehensive immigration reform.

``Congress is spending hundreds of billions of dollars to create new jobs, which should go to legal
U.S.~workers,'' Mr.~Shuler said. ``Americans have demanded that Congress do something about illegal
immigration, and enforcing existing laws is a good first step.''

White House officials do not rule out the possibility of an immigration overhaul before midterm
Congressional campaigns are in full swing next year. Some officials, however, say passage will more
likely come in 2011, when Mr.~Obama hopes to tap his broad support among Hispanic voters as he
begins his run for re-election.

On the other hand, Robert Menendez of New Jersey, the only Hispanic among Senate Democrats, says it
ought to be this year. ``I think it is one of those issues that if you don't pass this year, it
slips several years away,'' Mr.~Menendez said.

In the five months since Mr.~Obama took office, he has used his administrative authority to reverse
several Bush administration policies widely criticized as doing little to stem illegal immigration
while wreaking havoc on immigrant families already here.

At the Migration Policy Institute meeting Wednesday, John T.~Morton, assistant secretary of homeland
security, talked about some of those changes, including new guidelines that make employers, rather
than workers, the target of workplace raids, as well as expanded humanitarian-release rules to keep
parents detained on immigration charges from being separated from small children.

Immigration advocates say that administrative changes are not enough and that Hispanics, a crucial
voting bloc, are holding Mr.~Obama to his commitment to winning a comprehensive overhaul.

Representative Jason Altmire, a Democrat whose district includes the economically hard-hit
Pittsburgh suburbs, says his constituents are holding him to a different commitment.

``By definition,'' Mr.~Altmire said, ``illegal immigrants are people who broke the law to get here.
So any effort to reward them with legal status and work permits would not be supported by the people
who elected me.''

\section{At Air-Conditioned News Conference, Obama Feels the Heat}

WASHINGTON -- The Rose Garden is surely one of the most picturesque places in the capital, so it may
have been seen as regrettable when the summertime heat prompted the White House to move President
Obama's first formal outdoor news conference into the air conditioning.

As it turned out, it may have been cooler outside, even with the humidity.

In five months on the job, Mr.~Obama has rarely experienced as combative and contentious an hour on
live television as he did on Tuesday afternoon. He delivered the news he came to deliver -- a
message on Iran, energy legislation and health care -- before his question-and-answer session
quickly became infused with an air of testiness.

``I know everybody here is on a 24-hour news cycle,'' Mr.~Obama snapped at one point, after being
asked repeatedly by Chuck Todd of NBC News why he had not spelled out any consequences to the
Iranian government in the wake of the violent protests there. ``I'm not, O.K?''

For a White House, the press briefing room may be one of the least desirable places to hold a
presidential briefing. The space is neither formal nor imposing. And the man at the podium is seldom
so interruptible, with reporters sitting almost close enough to touch him.

The room may only be a 30-second walk from the Oval Office, but it is the equivalent of an away game
for the president. The small and crowded area -- normally home to the daily press briefing -- is the
only place in the West Wing where reporters play on their home turf, where the names of news
organizations are etched on each seat.

The close quarters produced more give-and-take and featured less decorum than the chandelier-lit
East Room, where Mr.~Obama's first three news conferences have been, or the sprawling Rose Garden,
where Tuesday's event was initially set to take place.

``What, are you the ombudsman for the White House press corps?'' Mr.~Obama said at one point,
chiding Jake Tapper of ABC News for trying to ask about two things, including following up on a
health care question the president had not directly answered.

As the reporter's question lingered, the president interjected: ``I got it. You're pitching. I'm
catching. I got the question.''

Later, as Mr.~Obama talked about the virtues of freedom of expression in Iran, Helen Thomas piped up
from her seat in the front row, ``Then why won't you allow the photos from the guards?'' (She
appeared to be speaking of his refusal to release photographs of abuse of prisoners by military
guards.)

``Hold on a second, Helen,'' said Mr.~Obama, who clearly had not planned to call on the
longest-serving scribe in the room. ``That's a different question.''

By the time one of Mr.~Obama's least-favorite topics was broached -- his occasional smoking habit --
the tone was set. It was a question the White House had prepared for, considering the president
signed a historic anti-smoking law one day earlier. But his smile did little to mask his agitation.

``You think it's neat to ask me about my smoking, as opposed to it being relevant to my new law,''
Mr.~Obama responded to his questioner, Margaret Talev of McClatchy Newspapers. ``That's fine. I
understand. It's an interesting human-interest story.''

As the president spoke, four of his top advisers watched from their seats on the side of the room.
There were smiles and occasional grimaces, with Rahm Emanuel, the chief of staff, smirking during
the cigarette question. The space is so tight that most White House aides have to tune in like
everyone else -- on their television sets.

It was Mr.~Obama's first daytime press conference, which Dana Milbank in the Washington Post
described as ``a new daytime drama'' that preempted soap operas on all three broadcast networks. He
dismissed the program as ``pre-packaged entertainment,'' particularly an exchange with a reporter
from a liberal Internet site who was invited in advance to relay a question submitted from a reader
in Iran. (Liberal bloggers long expressed suspicions that the Bush White House orchestrated news
conference questions in advance, which it did not. On Tuesday, there was no outcry.)

``Since we're on Iran,'' Mr.~Obama said, ``I know Nico Pitney is here from the Huffington Post.''

``Thank you, Mr.~President,'' replied Mr.~Pitney, who has been blogging about the violent protests
in Iran following the disputed presidential elections. ``I wanted to use this opportunity to ask you
a question directly from an Iranian.''

In recent administrations, presidents have only occasionally convened press conferences in the
briefing room, where the raucous setting can change the vibe -- and the news -- that ultimately
comes out.

``The briefing room is small -- it's really a studio -- so everything in its close quarters feels
more contentious and combative,'' said Ari Fleischer, who was a press secretary in President George
W.~Bush's first term. ``In contrast, the East Room feels stately no matter how tough the
questions.''

``But no matter what,'' he added, ``the setting can't save the president when events are bad or
getting worse.''

To several officials inside the Obama White House, the news conference was a success. It produced
front-page headlines denouncing the violent protests in Iran, which he condemned in his most
forceful terms yet. It also offered a more detailed explanation of what he would -- and would not --
accept in negotiations over his health care plan.

Still, the heat inside the building was noticeable, which is why it may be one of the last times the
president uses the briefing room for a formal press conference. Next time, it could well be back to
the Rose Garden.

\section{Strike Reportedly Missed Chief of Pakistani Taliban by Hours}

SLAMABAD, Pakistan -- A United States drone strike on a funeral in Pakistan's tribal areas missed
the leader of the Pakistani Taliban, Baitullah Mehsud, by hours on Tuesday, a Pakistani security
official said Wednesday.

Mr.~Mehsud was not present at the time of the attack, but had gone to pay his respects to a Taliban
commander killed in another American drone strike earlier the same day, the official said.

Though the strike on the funeral appeared to have included only two midlevel Taliban leaders among
the scores killed, it presented a clear blow to Mr.~Mehsud's operation, showing the deadly proximity
of the drone attacks to his areas and even the possibility that he was a target.

The Pakistani military has stepped up its operations against Mr.~Mehsud and his followers in South
Waziristan, mostly with airstrikes of its own.

A Pakistani police official speculated that the American drone attack could have been coordinated
with Pakistani officials, but could not confirm it. A Pakistani intelligence official, however, said
that Pakistan had been coordinating drone attacks with the Americans for several months.

In the past, most drone attacks were focused around Wana, the capital of South Waziristan, and
directed at targets that posed an immediate threat to the United States. These included foreign
members of Al Qaeda or Taliban commanders who helped coordinate cross-border attacks on American and
NATO forces in Afghanistan.

Recently, however, the drones have seemed to home in on Mr.~Mehsud and his groups, which have posed
a growing threat to the security of Pakistan through scores of suicide attacks in the country.

The apparent switch follows months of public and private complaints by Pakistani officials that the
Americans were interested in using the drones only against individuals who threatened their own
security interests in Afghanistan and beyond, not militants like Mr.~Mehsud who had struck inside
Pakistan.

That may be changing. The drones, which are operated remotely from the United States, carried out
two strikes in Mr.~Mehsud's area in South Waziristan on Tuesday.

The first strike hit a compound in Ladha, a village in the mountains of western Pakistan that is
Mr.~Mehsud's sanctuary. It killed Khwazh Wali, a Taliban leader who was a close aide to Mr.~Mehsud,
and three others.

Later Tuesday, Mr.~Wali's body was taken for burial to the village of Zangara, east of Makeen, where
people, apparently including Mr.~Mehsud, went to pay tribute to the Taliban fighter. It was unclear
how long after Mr.~Mehsud had left the village that the second drone sent three missiles hurtling
into the large crowd.

Estimates of the number of dead varied. One security official said that as many as 80 people had
been killed, but a local resident in a nearby town said the number was closer to 50. Pakistani
television reported that more than 100 had been killed.

A resident supportive of Mr.~Mehsud said by telephone from the area that none of the most senior
leaders of Mr.~Mehsud's group had been killed.

Though Pakistani television networks reported widely that top commanders, including Qari Hussein and
Sangeen Zadran, the Afghan commander close to Mr.~Mehsud, had been killed, the resident in the area
denied those men had been at the funeral.

Their deaths would have represented a serious blow to the Taliban, which relies on Mr.~Hussein's
expertise to train suicide bombers.

Instead, the highest ranking Taliban commander killed in the attack on the funeral appeared to a
leader known as Commander Bilal. He was in charge of Mr.~Mehsud's Shabi Khel tribe and a central figure in the Taliban movement, according to a Mehsud Taliban fighter who spoke by telephone on Wednesday evening.

\section{Gov. Sanford Admits Affair and Explains Disappearance}

COLUMBIA, S.C. -- Gov. Mark Sanford of South Carolina said Wednesday that he had been having an
extramarital affair with a woman in Argentina for the last year, ending the mystery surrounding his
disappearance over Father's Day weekend and considerably dampening his prospects for a national
political career.

But his confession and apology, in a rambling, nationally televised news conference, left other
mysteries unsolved, like whether he had lied to his staff members as late as Monday about his
whereabouts, whether the affair had definitively ended, whether he would resign from the
governorship and whether he would even have acknowledged the affair had he not been met at the
airport in Atlanta by a reporter upon his return.

Mr.~Sanford, 49, admitted that he had been in Buenos Aires since Thursday, not hiking on the
Appalachian Trail, as his staff members had said.

Standing in the rotunda of the South Carolina Statehouse, the governor, a Republican who had been
considered a possible presidential candidate in 2012, teared up as he spoke, taking more than seven
minutes to apologize before getting to the crux of the matter.

``The bottom line is this,'' he said. ``I have been unfaithful to my wife.''

The governor's wife, Jenny, 46, who did not attend the news conference, issued a statement later in
the day saying that while she loves her husband, she asked him to leave the family two weeks ago in
a trial separation, though she still believes the marriage can be repaired. The couple have four
sons, the youngest 10.

``We reached a point where I felt it was important to look my sons in the eyes and maintain my
dignity, self-respect, and my basic sense of right and wrong,'' she said. Because of the separation,
she said, she did not know where he was in the last week.

The governor, who raised his national profile by opposing the Obama administration's economic
stimulus plan, said he would resign from his position as chairman of the Republican Governors
Association. He will be succeeded by Gov. Haley Barbour of Mississippi. A reporter tried to ask
Mr.~Sanford if he would resign from the governor's office, but he did not answer.

Coming a week after the admission of an extramarital affair by Senator John Ensign, a Nevada
Republican who had also begun exploring a presidential run in 2012, the governor's acknowledgment
was yet another blow to Republican hopes for a strong field of challengers to President Obama.

``Personal circumstances over the course of the last week have managed to shrink the front line of
the 2012 possible contender list by 30 percent,'' said Phil Musser, a former executive director of
the Republican Governors Association.

Mr.~Sanford is regarded as a political lone wolf and has made numerous enemies even within his own
party, which controls both houses of the state legislature. Scott H.~Huffmon, a political science
professor at Winthrop University in Rock Hill, S.C., said it was unclear whether there would be
pressure on the governor to resign. ``His opponents,'' Dr.~Huffmon said, ``are sitting back trying
to figure out if they'd be better off with a completely emasculated governor to deal with or if
they'd be better off with Andr\'e Bauer,'' the lieutenant governor, who would take office if
Mr.~Sanford resigned.

Mr.~Sanford made at least one state-sponsored trip to Argentina during the period of his
relationship. In an interview late Wednesday, Daniel Scioli, the governor of Buenos Aires Province,
said he met with Mr.~Sanford on June 26, 2008, in La Plata, a town outside Buenos Aires. Mr.~Scioli
said the request for the meeting had come from Mr.~Sanford's office via the United States Embassy.

The governor was not known as a moralist but has frowned on infidelity and as a congressman voted to
impeach President Bill Clinton after the Monica Lewinsky affair. ``He lied under a different oath,
and that's the oath to his wife,'' Mr.~Sanford said at the time on CNN. ``So it's got to be taken
very, very seriously.''

Mr.~Sanford and his wife had joined an intensive Bible study group for couples in the last few
months, according to William H.~Jones, president of Columbia International University, a
conservative evangelical college. Dr.~Jones said the governor and his wife had been encouraged to
join the Bible study by a longtime friend, Cubby Culbertson, a businessman in Columbia who teaches
the Bible and who was thanked by the governor in his news conference.

At the news conference, Mr.~Sanford said his friendship with the unnamed woman began eight years ago
and became a romance about a year ago. He said he had seen her three times since then. The
relationship was ``discovered'' five months ago, he said, and he had been trying to reconcile with
his wife.

Mr.~Sanford strongly implied that he had ended the affair. ``The one thing that you really find is
that you absolutely want resolution,'' he said. ``And so oddly enough, I spent the last five days of
my life crying in Argentina.''

On Wednesday afternoon, The State, the leading newspaper in Columbia, published on its Web site
several e-mail messages it said it obtained in December between Mr.~Sanford and a Buenos Aires woman
the newspaper identified only as Maria.

In one of the messages, the governor describes himself as being in a ``hopelessly impossible
situation of love'' and stops just short of going into what he describes as sexual details of their
encounters.

After a barrage of news media requests about the missing governor began Monday, the governor's
spokesman, Joel Sawyer, released a statement on Monday afternoon saying that the governor was taking
some time to recharge after the stimulus battle and to work on ``a couple of projects that have
fallen by the wayside.'' Ms.~Sanford told The Associated Press that her husband had gone somewhere
over the Father's Day weekend to write, but that she did not know where.

Mr.~Sawyer played down the controversy as a creation of Mr.~Sanford's political enemies.

Then, around 10 p.m. Monday, Mr.~Sawyer sent a ``high priority'' e-mail alert to reporters that
Mr.~Sanford was hiking on the 2,100-mile Appalachian Trail. Mr.~Sawyer flatly denied a television
news report that the governor had been seen boarding a plane at the airport in Atlanta.

At the news conference, the governor said his staff members had based their statements on tentative
scheduling information he had given before he left.

Mr.~Sanford's rivals immediately pounced on the apparent confusion in the governor's office.

``The people of this state deserve complete honesty from Governor Sanford,'' said State Senator John
C.~Land III, the Senate Democratic leader, in a statement issued Wednesday morning before the news
conference. ``Never in my 32 years as a state senator have I witnessed a governor and his staff act
in a more dishonest, secretive and bizarre manner.''

But Mr.~Land said that he doubted there would be pressure for Mr.~Sanford to go because a
resignation would mean that State Senator Glenn McConnell, the powerful president pro tem, would
have to become the lieutenant governor, a relatively powerless position.

Mr.~Sanford recently lost a high-profile battle to divert \$700 million in federal stimulus money
toward reducing the state deficit, challenging the Obama administration on the issue.

Critics said he was simply seeking to raise his national profile, but the governor maintained that
his primary goal was to strengthen the executive office in South Carolina, where the governor has
few powers.

Mr.~Sanford has long been known as an iconoclast. As a congressman, he slept on a futon in his
office. To showcase his opposition to pork-barrel spending, he once brought two live piglets onto
the floor of the legislature.

Still, many were shocked by his announcement. ``I never figured Sanford for anything like this, said
Neal D.~Thigpen, a political science professor at Francis Marion University in Florence, S.C.

\section{U.S. Objects to China's Web Filtering}

The Obama administration lodged a formal protest on Wednesday with the Chinese government over its
plan to force all computers sold in China to come with software that blocks access to certain Web
sites.

Commerce Secretary Gary Locke and Ron Kirk, the trade representative, sent a letter to officials in
two Chinese ministries asking them to rescind a rule about the software that is set to take effect
on July 1.

Chinese officials have said that the filtering software, known as Green Dam-Youth Escort, is meant
to block pornography and other ``unhealthy information.''

In part, the American officials' complaint framed this as a trade issue, objecting to the burden put
on computer makers to install the software with little notice. But it also raised broader questions
about whether the software would lead to more censorship of the Internet in China and restrict
freedom of expression.

``China is putting companies in an untenable position by requiring them, with virtually no public
notice, to pre-install software that appears to have broad-based censorship implications and network
security issues,'' Mr.~Locke said in a news release. The government did not release the text of the
letter.

The letter, by two cabinet-rank officials, represents an escalation of the concern over the software
plan. Last week United States officials met with their Chinese counterparts in Beijing to raise
objections to the new policy.

Ed Black, president of the Computer and Communications Industry Association, one of several trade
groups that have objected to the Chinese plan, said the letter represented a significant change in
American policy.

``The issue of Internet freedom and openness was something that should have been at the top of the
U.S.~international agenda and hasn't been,'' Mr.~Black said. ``This administration is far more in
tune with and ready to support Internet openness.''

China already has an elaborate system that blocks access to sites that discuss delicate topics like
the Dalai Lama and Falun Gong, the banned spiritual movement.

In their statement, the American officials rejected the argument that the software was simply a way
to block pornography.

``Protecting children from inappropriate content is a legitimate objective, but this is an
inappropriate means and is likely to have a broader scope,'' Mr.~Kirk said in the statement.
``Mandating technically flawed Green Dam software and denying manufacturers and consumers freedom to
select filtering software is an unnecessary and unjustified means to achieve that objective.''

Security experts have expressed concerns that once installed, the software might also be used to
block other sorts of content or even to monitor the online activities of citizens.

The letter suggested that China's move might violate World Trade Organization rules because American
companies were given only six weeks' notice to comply. While formal complaints to the trade
organization are difficult and cumbersome, pointing to the regulations is another signal that the
United States will continue to pursue the issue.

With only one week before the new rules are to go into effect, it is unclear if American computer
companies will comply.

Pamela Bonney, a spokeswoman for Hewlett-Packard, said the company was still studying the rules and
seeking clarification. A spokesman for Dell did not return calls seeking comment.

Separately, access to Google's main search engine at Google.com and other services like Gmail was
temporarily blocked in China on Wednesday. It was restored a few hours later.

Access to foreign Web sites in China can be erratic, and determining whether the government is
responsible can be difficult.

It is not clear whether the blocking of Google's sites is related to a dispute that erupted last
week between Google and Chinese authorities. The Chinese government disabled some search functions
of Google's Chinese-language search engine, Google.cn, saying the site offered too many links to
pornographic material.

Google's license to operate in China requires that it not show pornographic sites.

\section{Iran Stepping Up Effort to Quell Election Protest}

TEHRAN -- Iranian officials stepped up efforts to crush the remaining resistance to a disputed
presidential election on Wednesday, as security forces overwhelmed a small group of protesters with
brutal beatings, tear gas and gunshots in the air. Intelligence agents shut down an office of a
defeated presidential candidate, saying it was a ``headquarters for a psychological war.''

The nation's leadership cast anyone refusing to accept the results of the race as an enemy of the
state. Analysts suggested that the unyielding response showed that Iran's leaders, backed by the
supreme leader, Ayatollah Ali Khamenei, had lost patience and that Iran was now, more than ever, a
state guided not by clerics of the revolution but by a powerful military and security apparatus.

President Mahmoud Ahmadinejad has maintained a low profile, but evidence suggests that he has filled
security agencies with crucial allies.

``What has been going on since 2005 is the shift of the center of power from the clergy to the
Pasdaran,'' or the Islamic Revolutionary Guards Corps, said a political analyst with years of
experience in Iran who feared retribution if identified. ``In a way one could say that Iran is no
longer a theocracy, but a government headed by military chiefs.''

Security agents continued to fan out across the country, detaining former government officials,
journalists, activists, young people and old, anyone seen as siding with those who reject the
conclusion that Mr.~Ahmadinejad won a landslide against the leading opposition candidate, Mir
Hussein Moussavi.

The official Iranian news agency reported that intelligence and security agents in Tehran concluded
that a Moussavi campaign office was used for ``illegal gatherings, the promotion of unrest, and
efforts to undermine the country's security,'' leading to speculation that Mr.~Moussavi could be
arrested. The news agency reported that ``the plotters have been arrested.''

The government also stepped up its efforts to block independent news coverage of events all across
the country. The government banned foreign news media members from leaving their offices, suspended
all press credentials for the foreign press, arrested a freelance writer for The Washington Times,
continued to hold a reporter for Newsweek and forced other foreign journalists to leave the country.

That made it difficult to ascertain exactly what happened when several hundred protesters tried to
gather outside the Parliament building late in the afternoon. Witnesses said they were met by a huge
force of riot police officers and Basij vigilantes, some on motorcycles and some in pickup trucks,
armed with sticks and chains. Witnesses said people were trapped and beaten as they tried to flee
down side streets.

``It was not possible to wait and see what happened,'' said one witness who asked for anonymity out
of fear of arrest. ``At one point we saw several riot police in black clothes walk towards a group
of people who looked like passers-by. Suddenly they pulled out their batons and began hitting them
without warning.''

Witnesses reported scenes of chaos and fear where riot police officers outnumbered demonstrators by
about four to one. Many wore masks to conceal their identities. The Basijis stopped people to check
their cellphones for video or pictures of the unrest.

``I saw one group of about 100 people who began chanting 'Death to the dictator' on one of the side
streets,'' said another witness who insisted on not being identified for fear of arrest. ``The
Basijis attacked them and beat them really bad.'' Unconfirmed reports of bloodshed and at least one
death flooded the Internet.

The authorities said they were moving to impose order and secure the rule of law. ``I was insisting
and will insist on implementation of the law,'' Ayatollah Khamenei said on national television.
``That means we will not go one step beyond the law. Neither the system nor the people will yield to
pressure at any price.''

The pressure exposed deep cracks in the opposition, but also sparked signs of entrenched resistance.
Early in the day, Mohsen Rezai, a former head of the elite Revolutionary Guards and a presidential
candidate, withdrew his charges of election fraud, saying that it was in the best interest of the
country to drop the matter. His decision helped the government's effort to cast the opposition as at
war with the state.

Then Mr.~Moussavi, the defeated candidate who embodied the hopes of reformers, posted a notice on
his Web site of a late afternoon rally in front of the Parliament, but he distanced himself from the
action, saying it was not organized by the reform movement. It is not clear how far Mr.~Moussavi, a
former prime minister who is essentially an insider thrust into the role of opposition, would go to
defy the system. He has not been seen since Thursday. So as the crackdown infuriates protesters,
there is a greater gap with their ostensible leader, political analysts said.

``People in the street have been radicalized, and I do not believe that most of them would today
subscribe to Moussavi's avowed platform,'' said a political analyst with years of experience in
Iran.

Instead, Mr.~Moussavi's wife, Zahra Rahnavard, a former university dean, continued to rally the
opposition movement. She posted a message on another Web site associated with her husband calling on
the public to stand firm while criticizing the government for acting ``as if martial law has been
imposed on the streets.''

The head of Iran's intelligence service announced that would-be terrorists, including two foreign
reporters, were detained all over the country planning explosions and acts of sabotage, the official
Iran news agency reported.

Iran's defense minister, Brig. Gen. Mostafa Mohammad Najjar, pointed to recent military maneuvers in
the Persian Gulf as proof that Iran could crush any foreign threat. The government's harsh response,
including the killing of at least 17 protesters, led to divisions among some conservatives who
criticized the armed attacks on unarmed civilians. The mayor of Tehran, Mohammad Baqer Qalibaf, a
former commander of the Iranian police, called on the government to authorize peaceful opposition
rallies. And the speaker of Parliament, Ali Larijani, a longtime conservative, accused the Guardian
Council, responsible for monitoring the elections, of bias and said most Iranians were suspicious of
the election results.

Instead of heeding calls for moderation, the government has conducted one of the harshest crackdowns
in its history. Dozens of former high-ranking officials have been jailed. The International Campaign
for Human Rights in Iran reported Wednesday that about 240 people, including 102 political figures,
were in jail. The government has said that it arrested 627 more people since the protests broke out.

Those arrested include officials who served from the founding of the Islamic republic in 1979, until
Mr.~Ahmadinejad's election in 2005: Behzad Nabavi, a former deputy speaker of Parliament; Mohsen
Aminzadeh, a key figure at the Intelligence Ministry for many years; Mostafa Tajzadeh, a deputy
interior minister during the presidency of Mohammad Khatami; Mohammad Ali Abtahi, a vice president
under Mr.~Khatami; and Abdullah Ramezanzadeh, Mr.~Khatami's spokesman. They were all close to
Mr.~Khatami, then threw their support behind Mr.~Moussavi.

\section{A Star Idolized and Haunted, Michael Jackson Dies at 50}

LOS ANGELES -- For his legions of fans, he was the Peter Pan of pop music: the little boy who
refused to grow up. But on the verge of another attempted comeback, he is suddenly gone, this time
for good.

Michael Jackson, whose quintessentially American tale of celebrity and excess took him from musical
boy wonder to global pop superstar to sad figure haunted by lawsuits, paparazzi and failed plastic
surgery, was pronounced dead on Thursday afternoon at U.C.L.A. Medical Center after arriving in a
coma, a city official said. Mr.~Jackson was 50, having spent 40 of those years in the public eye he
loved.

The singer was rushed to the hospital, a six-minute drive from the rented Bel-Air home in which he
was living, shortly after noon by paramedics for the Los Angeles Fire Department. A hospital
spokesman would not confirm reports of cardiac arrest. He was pronounced dead at 2:26 pm.

As with Elvis Presley or the Beatles, it is impossible to calculate the full effect Mr.~Jackson had
on the world of music. At the height of his career, he was indisputably the biggest star in the
world; he has sold more than 750 million albums. Radio stations across the country reacted to his
death with marathon sessions of his songs. MTV, which grew successful in part as a result of
Mr.~Jackson's groundbreaking videos, reprised its early days as a music channel by showing his
biggest hits.

From his days as the youngest brother in the Jackson 5 to his solo career in the 1980s and early
1990s, Mr.~Jackson was responsible for a string of hits like ``I Want You Back,'' ``I'll Be There''
``Don't Stop 'Til You Get Enough'' ``Billie Jean'' and ``Black or White'' that exploited his high
voice, infectious energy and ear for irresistible hooks.

As a solo performer, Mr.~Jackson ushered in the age of pop as a global product -- not to mention an
age of spectacle and pop culture celebrity. He became more character than singer: his sequined
glove, his whitened face, his moonwalk dance move became embedded in the cultural firmament.

His entertainment career hit high-water marks with the release of ``Thriller,'' from 1982, which has
been certified 28 times platinum by the Recording Industry Association of America, and with the
``Victory'' world tour that reunited him with his brothers in 1984.

But soon afterward, his career started a bizarre disintegration. His darkest moment undoubtedly came
in 2003, when he was indicted on child molesting charges. A young cancer patient claimed the singer
had befriended him and then groped him at his Neverland estate near Santa Barbara, Calif., but
Mr.~Jackson was acquitted on all charges.

Reaction to his death started trickling in from the entertainment community late Thursday.

``I am absolutely devastated at this tragic and unexpected news,'' the music producer Quincy Jones
said in a statement. ``I've lost my little brother today, and part of my soul has gone with him.''

Berry Gordy, the Motown founder who helped develop the Jackson 5, told CNN that Mr.~Jackson, as a
boy, ``always wanted to be the best, and he was willing to work as hard as it took to be that. And
we could all see that he was a winner at that age.

Tommy Mottola, a former head of Sony Music, called Mr.~Jackson ``the cornerstone to the entire music
business.''

``He bridged the gap between rhythm and blues and pop music and made it into a global culture,''
said Mr.~Mottola, who worked with Mr.~Jackson until the singer cut his ties with Sony in 2001.

Impromptu vigils broke out around the world, from Portland, Ore., where fans organized a one-gloved
bike ride (``glittery costumes strongly encouraged'') to Hong Kong, where fans gathered with candles
and sang his songs.

In Los Angeles, hundreds of fans -- some chanting Mr.~Jackson's name, some doing the ``Thriller''
dance -- descended on the hospital and on the hillside house where he was staying.

Jeremy Vargas, 38, hoisted his wife, Erica Renaud, 38, on his shoulders and they danced and bopped
to ``Man in the Mirror'' playing from an onlooker's iPod connected to external speakers -- the boom
boxes of Mr.~Jackson's heyday long past their day.

``I am in shock and awe,'' said Ms.~Renaud, who was visiting from Red Hook, Brooklyn, with her
family. ``He was like a family member to me.''

Dreams of a Comeback

Mr.~Jackson was an object of fascination for the news media since the Jackson 5's first hit, ``I
Want You Back,'' in 1969. His public image wavered between that of the musical naif, who wanted only
to recapture his youth by riding on roller-coasters and having sleepovers with his friends, to the
calculated mogul who carefully constructed his persona around his often-baffling public behavior.

Mr.~Jackson had been scheduled to perform 50 concerts at the O2 arena in London beginning next month
and continuing into 2010. The shows, which quickly sold out, were positioned as a comeback, with the
potential to earn him up to \$50 million, according to some reports.

But there had also been worry and speculation that Mr.~Jackson was not physically ready for such an
arduous run of concerts, and his postponement of the first of those shows to July 13 from July 8
fueled new rounds of gossip about his health. Nevertheless, he was rehearsing Wednesday night at the
Staples Center in downtown Los Angeles. ``The primary reason for the concerts wasn't so much that he
was wanting to generate money as much as it was that he wanted to perform for his kids,'' said
J.~Randy Taraborrelli, whose biography, ``Michael Jackson: The Magic and the Madness,'' was first
published by Citadel in 1991. ``They had never seen him perform before.''

Mr.~Jackson's brothers, Jackie, Tito, Jermaine, Marlon and Randy, have all had performing careers,
with varying success, since they stopped performing together. (Randy, the youngest, replaced
Jermaine when the Jackson 5 left Motown.) His sisters, Rebbie, La Toya and Janet, are also singers,
and Janet Jackson has been a major star in her own right for two decades. They all survive him, as
do his parents, Joseph and Katherine Jackson, of Las Vegas, and three children: Michael Joseph
Jackson Jr., Paris Michael Katherine Jackson, born to Mr.~Jackson's second wife, Deborah Jeanne
Rowe, and Prince Michael Jackson II, the son of a surrogate mother. Mr.~Jackson was also briefly
married to Lisa Marie Presley, the daughter of Elvis Presley.

A spokesman for the Los Angeles Police Department said the department assigned its robbery and
homicide division to investigate the death, but the spokesman said that was because of Mr.~Jackson's
celebrity.

``Don't read into anything,'' the spokesman told reporters gathered outside the Bel-Air house. He
said the coroner had taken possession of the body and would conduct an investigation.

At a news conference at the hospital, Jermaine Jackson spoke to reporters about his brother. ``It is
believed he suffered cardiac arrest at his home,'' he said softly. A personal physician first tried
to resuscitate Michael Jackson at his home before paramedics arrived. A team of doctors then tried
to resuscitate him for more than an hour, his brother said.

``May Allah be with you always,'' Jermaine Jackson concluded, his gaze aloft.

In Gary, Ind., hundreds of people descended upon the squat clapboard house were Mr.~Jackson spent
his earliest years. There were tears, loud wails, and quiet prayers as old neighbors joined hands
with people who had driven in from Chicago and other nearby towns to pay their respects.

``Just continue to glorify the man, Lord,'' said Ida Boyd-King, a local pastor who led the crowd in
prayer. ``Let's give God praise for Michael.''

Shelletta Hinton, 40, drove to Gary from Chicago with her two young children. She said they had met
Mr.~Jackson in Gary a couple of years ago when he received a key to the city. ``We felt like we were
close to Michael,'' she said. ``This is a sad day.''

As dusk set in, mourners lighted candles and placed them on the concrete doorstep. Some left teddy
bears and personal notes. Doris Darrington, 77, said she remembered seeing the Jackson 5 so many
times around Gary that she got sick of them. But she, too, was feeling hurt by the sudden news of
Mr.~Jackson's death.

``He has always been a source of pride for Gary, even though he wasn't around much,'' she said.
``The older person, that's not the Michael we knew. We knew the little bitty boy with the big Afro
and the brown skin. That's how I'll always remember Michael.''

Michael Joseph Jackson was born in Gary on Aug.~29, 1958. The second youngest of six brothers, he
began performing professionally with four of them at the age of 5 in a group that their father, Joe,
a steelworker, had organized the previous year. In 1968, the group, originally called the Jackson
Brothers, was signed by Motown Records. The Jackson 5 was an instant phenomenon. The group's first
four singles -- ``I Want You Back,'' ``ABC,'' ``The Love You Save'' and ``I'll Be There'' -- all
reached No.~1 on the pop charts in 1970, a feat no group had accomplished before. And young Michael
was the center of attention: he handled virtually all the lead vocals, danced with energy and
finesse, and displayed a degree of showmanship rare in a performer of any age.

In 1971, Mr.~Jackson began recording under his own name, while continuing to perform with his
brothers. His recording of ``Ben,'' the title song from a movie about a boy and his homicidal pet
rat, was a No.~1 hit in 1972.

The brothers (minus Michael's older brother Jermaine, who was married to the daughter of Berry
Gordy, Motown's founder and chief executive) left Motown in 1975 and, rechristened the Jacksons,
signed to Epic, a unit of CBS Records. Three years later, Michael made his movie debut as the
Scarecrow in the screen version of the hit Broadway musical ``The Wiz.'' But movie stardom proved
not to be his destiny.

A Solo Sensation

Music stardom on an unprecedented level, however, was. Mr.~Jackson's first solo album for Epic,
``Off the Wall,'' released in 1979, yielded two No.~1 singles and sold seven million copies, but it
was a mere prologue to what came next. His follow-up, ``Thriller,'' released in 1982, became the
best-selling album of all time and helped usher in the music video age. The video for title track,
directed by John Landis, was an elaborate horror-movie pastiche that was more of a mini-movie than a
promotional clip.

Seven of the nine tracks on ``Thriller'' were released as singles and reached the Top 10. The album
spent two years on the Billboard album chart and sold an estimated 100 million copies worldwide. It
also won eight Grammy Awards.

The choreographer and director Vincent Paterson, who directed Mr.~Jackson in several videos,
recalled watching him rehearse a dance sequence for four hours in front of a mirror until it felt
like second nature.

``That's how he developed the moonwalk, working on it for days if not weeks until it was organic,''
he said. ``He took an idea that he had seen some street kids doing and perfected it.''

Mr.~Jackson's next album, ``Bad,'' released in 1987, sold eight million copies and produced five
No.~1 singles and another state-of-the-art video, this one directed by Martin Scorsese. It was a
huge hit by almost anyone else's standards, but an inevitable letdown after ``Thriller.''

It was at this point that Mr.~Jackson's bizarre private life began to overshadow his music. He would
go on to release several more albums and, from time to time, to stage elaborate concert tours. And
he would never be too far from the public eye. But it would never again be his music that kept him
there.

Even with the millions Mr.~Jackson earned, his eccentric lifestyle took a severe financial toll. In
1988 Mr.~Jackson paid about \$17 million for a 2,600-acre ranch in Los Olivos, Calif., 125 miles
northwest of Los Angeles. Calling it Neverland after the mythical island of Peter Pan, he outfitted
the property with amusement-park rides, a zoo and a 50-seat theater, at a cost of \$35 million,
according to reports, and the ranch became his sanctum.

But Neverland, and Mr.~Jackson's lifestyle, were expensive to maintain. A forensic accountant who
testified at Mr.~Jackson's molesting trial in 2005 said Mr.~Jackson's annual budget in 1999 included
\$7.5 million for personal expenses and \$5 million to maintain Neverland. By at least the late
1990s, he began to take out huge loans to support himself and pay debts. In 1998, he took out a loan
for \$140 million from Bank of America, which two years later was increased to \$200 million.
Further loans of hundreds of millions followed.

The collateral for the loans was Mr.~Jackson's 50 percent share in Sony/ATV Music Publishing, a
portfolio of thousands of songs, including rights to 259 songs by John Lennon and Paul McCartney,
considered some of the most valuable properties in music.

In 1985, Mr.~Jackson paid \$47.5 million for ATV, which included the Beatles songs -- a move that
estranged him from Mr.~McCartney, who had advised him to invest in music rights -- and 10 years
later, Mr.~Jackson sold 50 percent of his interest to Sony for \$90 million, creating a joint
venture, Sony/ATV. Estimates of the catalog's value exceed \$1 billion.

Last year, Neverland narrowly escaped foreclosure after Mr.~Jackson defaulted on \$24.5 million he
owed on the property. A Los Angeles real estate investment company, Colony Capital L.L.C., bought
the note, and put the title for the property into a joint venture with Mr.~Jackson.

A Scandal's Heavy Toll

In many ways, Mr.~Jackson never recovered from the child molesting trial, a lurid affair that
attracted media from around the world to watch as Mr.~Jackson, wearing a different costume each day,
appeared in a small courtroom in Santa Maria, Calif., to listen as a parade of witnesses spun a
sometimes-incredible tale.

The case ultimately turned on the credibility of Mr.~Jackson's accuser, a 15-year-old cancer
survivor who said the defendant had gotten him drunk and molested him several times. The boy's
younger brother testified that he had seen Mr.~Jackson groping his brother on two other occasions.

After 14 weeks of such testimony and seven days of deliberations, the jury returned not-guilty
verdicts on all 14 counts against Mr.~Jackson: four charges of child molesting, one charge of
attempted child molesting, one conspiracy charge and eight possible counts of providing alcohol to
minors. Conviction could have brought Mr.~Jackson 20 years in prison. Instead, he walked away a free
man to try to reclaim a career that at the time had already been in decline for years.

After his trial, Mr.~Jackson largely left the United States for Bahrain, the island nation in the
Persian Gulf, where he was the guest of Sheik Abdullah, a son of the ruler of the country, King
Hamad bin Isa al-Khalifa. Mr.~Jackson would never return to live at his ranch. Instead he remained
in Bahrain, Dubai and Ireland for the next several years, managing his increasingly unstable
finances. He remained an avid shopper, however, and was spotted at shopping malls in the black robes
and veils traditionally worn by Bahraini women.

Despite the public relations blow of his trial, Mr.~Jackson and his ever-changing retinue of
managers, lawyers and advisers never stopped plotting his return.

By early this year, Mr.~Jackson was living in a \$100,000-a-month mansion in Bel-Air, to be closer
to ``where all the action is'' in the entertainment business, his manager at the time, Tohme Tohme,
told The Los Angeles Times. He was also preparing for his upcoming London shows.

''He was just so excited about having an opportunity to come back,'' said Mr.~Paterson, the director
and choreographer.

Despite his troubles, the press and the public never abandoned the star. A crowd of paparazzi and
onlookers lined the street outside Mr.~Jackson's home as the ambulance took him to the hospital.

\textit{Reporting was contributed by John M.~Broder from Washington; Randal C.~Archibold from Los Angeles;
Susan Saulny from Gary, Ind.; and Melena Ryzik, Ben Sisario, Brian Stelter and Peter Keepnews from
New York.}

\texttt{This article has been revised to reflect the following correction:}

\texttt{Correction: June 27, 2009 An article on Friday about the death of Michael Jackson misstated the
number of songs from his album ``Off the Wall'' that became No.~1 singles. There were two, not four.
The article also misstated part of a comment that Mr.~Jackson's brother Jermaine offered for
Mr.~Jackson after speaking with reporters. He said, ``May Allah be with you always,'' not ``May our
love be with you always.''}

\section{Young Japanese Raise Their Voices Over Economy}

TOKYO -- A group of young people recently gathered in a darkened park here. Holding placards and
megaphones, they chanted slogans condemning the Japanese government and a lack of jobs and
opportunity.

The scene, which is repeated often in the gritty Tokyo neighborhood of Koenji, is nothing close to
the protests that have recently shaken Iran. Indeed, the protests would hardly raise an eyebrow in
most parts of the world, but in this country, which values conformity, they represent a stark
departure from the norm. Since the 1960s, when youth protests turned violent, even the mildest form
of protests by young people has been viewed as taboo.

But the pain of recession is changing that, giving rise to a new activism among Japan's youth, who
have long been considered apathetic.

``I'm here because I want to change society,'' one leader, Yoshihiro Sato, 28, recently shouted to a
crowd of about 50. ``Will you join me?''

Unlike the '60s generation, which agitated to change the bourgeois basis of Japanese society,
Mr.~Sato and other young people are today fighting to join it. They are demanding greater
professional opportunities, more job security and a stronger social safety net.

After so many decades without a grass-roots movement, protests are so rare here that many who wish
to take part require basic training.

The Tokyo-based Pacific Asia Research Center, an institute that typically runs seminars on social
issues like poverty, organized the recent march. After a surge of interest from young people who
said they wanted to get more involved in social issues but did not know how, the center started
offering what it says is Japan's first activist training program. The sessions include poster-making
and campaigning on the Web.

``Once we're done, we'll overrun Japan with demonstrations,'' Seiko Uchida, the head of a research
center, told a cheering crowd.

That may be hyperbole, but the deteriorating economy has inarguably affected young people more than
any other demographic. Unemployment was 9.6 percent in April for Japanese aged 15 to 24, compared
with 5 percent unemployment over all.

But unemployment and welfare benefits are sparse in Japan. And government spending is skewed toward
pensions and health care for older voters rather than programs that might train young workers or
help them support their families.

In the first quarter of the year, Japan's economy shrank a devastating 14.2 percent on an annual
basis, as exports slumped because of the global economic slowdown. Many of those who lost their jobs
were younger people in precarious ``temporary'' positions that were the product of a decade-long
deregulation.

The disparity has fueled generational friction, particularly between those who reaped the fruits of
Japan's rapid postwar growth and younger Japanese who came of age in Japan's ``Lost Decade'' of the
1990s, when the country's economic growth stagnated, and during its anemic recovery.

When companies like Canon and Toyota Motor started to fire temporary factory workers late last year,
a handful of the workers lashed out publicly, confronting managers at factory gates, often in front
of TV cameras. Others brought a flurry of lawsuits against former employers.

Over the New Year's holidays, about 500 laid-off temporary workers who had lost their homes
congregated at a park in the center of Tokyo, building an impromptu tent city next to the offices of
the Labor Ministry. The scene led to a media frenzy and national soul-searching on the plight of
young Japanese.

Today, workers are rushing to unionize, and the Japanese Communist Party says it is getting about
1,000 new members a month, many of them disgruntled young people.

Masahiro Mukasa, a struggling techno musician, started a union for freelance artists and musicians
in December. The Indy Union intends to help members negotiate with particularly abusive employers.

``People think musicians just have a good time. But we need to make a living, too,'' Mr.~Mukasa
said. ``I want to show that our livelihoods are at stake in this bad economy.''

Still, in a society that values conformity and order, most Japanese remain deeply averse to
confrontation and protest, and there is nothing approaching what could be called a mass movement.

``Japanese feel it's shameful to get involved in protests,'' said Makoto Yuasa, a longtime activist.
``Many still look at us suspiciously, like we might be making bombs.''

Still, ``this is the most significant rise in activism I've seen in years,'' said Yoshitaka Mouri, a
professor at the Tokyo University of the Arts, who has been following the rising protest culture.
``A movement is brewing among young Japanese.''

Hajime Matsumoto, an activist who operates from a thrift shop in Koenji, has amassed a large
following at his protests and rallies. Some like-minded Japanese have opened their own stores
alongside Mr.~Matsumoto's, huddling after hours to help hatch protest plans, turning Koenji into a
center of activism.

``The poor man's revolt has finally begun!'' Mr.~Matsumoto shouted at a recent demonstration,
banging on a full set of drums perched atop a piece of plywood on wheels. His message: even poor
people deserve a good life. ``If we all get together, we can bring about change!''

Some experts question how much political influence demonstrators will wield. Few expect them to be a
big force at the ballot box later this year when Japan holds parliamentary elections.

Young people are outnumbered by older voters, and are concentrated in cities, where ballots carry
less weight, proportionally, than in the sparsely populated countryside.

Still, the nascent focus on worker and generational rights is a break from the years under the
former nationalist prime ministers Shinzo Abe and Junichiro Koizumi. Then, stoked by nationalist
rhetoric from politicians and government officials, youth seemed to swerve to the right.

To quell the rising anger, the government has increased spending on programs for younger people. In
his latest economic stimulus package, Prime Minister Taro Aso pledged 1.9 trillion yen (\$19
billion) in programs to raise youth employment. He is also prodding companies to elevate temporary
workers to full-time status. Such actions have done little to change the economic issues.

For now, the public has some sympathy with the protesters, and the rallies get heavy coverage in
local media.

``I support these young people,'' said Masaaki Saito, 60, an owner of a small electronics store in
Koenji who took part in student protests in his youth.

``It's been a long time,'' he said, ``but Japan's youth are getting their voice back.''

\section{Telecom Firms Back Standard Phone Charger in Europe}

BRUSSELS (Reuters) - Top mobile telephone suppliers have agreed to back an EU-wide harmonization of
phone chargers, the European Commission said on Monday, hailing the pact as good news for consumers
and the environment.

The agreement by Nokia, Sony Ericsson and other industry majors will mean phones compatible with
standard charging devices are available in Europe from next year, said the EU executive, which has
pushed for such a deal.

``People will not have to throw away their charger whenever they buy a new phone,'' said EU Industry
Commissioner Guenter Verheugen, estimating that unwanted phone accessories accounted for thousands
of tons of waste in Europe each year.

The Commission said the agreement would involve the creation of an EU norm, and that the new
generation of mobile phones would use a standard micro-USB socket to ensure compatibility.

There are an estimated 400 million mobile phones in Europe, with 185 million bought each year.

The chargers will be usable only for data-enabled phones, which have more capability than just
standard calls and SMS texts. Data-enabled phones are expected to account for almost half of all new
mobile handset purchases in 2010.

The Commission hopes that as people discard their old handsets, within three to four years all
data-enabled phones in Europe will be using standardized chargers.

New data-enabled phones will come with a standardized charger but after an unspecified time the two
items will be sold separately, industry group DigitalEurope said.

The director of DigitalEurope, Bridget Cosgrave, said the price of such chargers had not yet been
determined.

Motorola, Apple, LG, NEC, Qualcomm, Research in Motion, Samsung and Texas Instruments have also
signed the agreement, the Commission said.

The 10 companies involved control 90 percent of the European market.

The standardized chargers will be compatible only with European phones, but Verheugen said he hoped
other countries would follow Europe's lead.

``We're assuming this new European initiative will have a knock-on effect globally and manufacturers
won't just be doing this on the European market,'' Verheugen said.

\section{Obama Seeks Popular Support for Health Reform}

WASHINGTON -- With Democrats deeply divided over health legislation, President Obama is trying to
enlist the nation's governors and his own army of grass-roots supporters in a bid to increase
pressure on lawmakers without getting himself mired in the messy battle playing out on Capitol Hill.

In a meeting last week with five governors -- including Republicans who may be more sympathetic to
health legislation than those on Capitol Hill -- Mr.~Obama privately urged them to serve as his
emissaries to Congress. He even coached them on the language they should use with lawmakers, two of
the governors said, advising them to avoid terms like ``rationing'' and ``managed care,'' which
evoke bitter memories of the Clintons' ill-fated health initiative.

The hourlong session in the Roosevelt Room was part of an intensifying but potentially risky White
House strategy to shift the health care debate away from Washington and to the states. On Wednesday,
Mr.~Obama will travel to Virginia to hold a town-hall-style meeting on health care -- his second in
two weeks -- that will include questions from online communities like Facebook and Twitter.

With members of Congress back in their districts for the Fourth of July recess, Mr.~Obama's
political group, Organizing for America, has recruited thousands of supporters to participate in
blood drives, raise money for medical research and volunteer at community health clinics this week,
all with the intent of sending reminders to lawmakers that the public wants action on health care.

``The main thing,'' David Axelrod, Mr.~Obama's senior adviser, said, ``is to involve as many people
as possible and demonstrate in a variety of ways the level and degree of intensity of support that
this has.'' Of Mr.~Obama's supporters, Mr.~Axelrod said, ``There's no issue that motivates them more
than health care.''

While this outside-the-Beltway strategy lets Mr.~Obama stay out of Democrats' internal fights -- for
now at least -- there are risks. If Mr.~Obama waits too long to exert his presidential muscle to
forge consensus on Capitol Hill, his moment of opportunity could pass. He could also lose control of
the final outcome if lawmakers cut backroom deals he dislikes, for example, by deciding to pay for
the expansion by taxing employee health benefits, a move that worries Mr.~Obama's political advisers
because it could cause the president to break a campaign promise.

Some Democrats are privately pushing the president to do more to bring his party in line. When Rahm
Emanuel, the White House chief of staff, went to Capitol Hill last week, the majority leader,
Senator Harry Reid of Nevada, pressed for the president to intervene more directly to settle
Democrats' disputes over Mr.~Obama's call for a government-run insurance plan to compete with the
private sector, two people familiar with the session said.

Mr.~Emanuel, in an e-mail message, acknowledged that some Democrats ``wanted more direct and
specific involvement,'' but said others were happy with the president's level of engagement, adding,
``We received a lot of advice.''

Over the last several weeks, Mr.~Obama has steadily increased his contact with lawmakers on health
care, even as he steers clear of specific policy disputes. He met privately with Senator Ron Wyden,
Democrat of Oregon, and telephoned Democrats on the Senate Finance Committee to check on their
progress and urge them to stick to his timetable for a final bill to reach his desk in October.

John D.~Podesta, who ran Mr.~Obama's transition to the presidency and consults closely with the
administration on the health bill, predicted that the White House would resist the urge to ``knock
heads and hammer consensus'' at least until after the Finance Committee produced a bill, sometime
after the Fourth of July holiday.

But if the panel, widely regarded as Mr.~Obama's best hope for a bipartisan measure, gets stuck or
further delayed, Mr.~Podesta and other Democrats say, Mr.~Obama will have to step in to broker a
deal.

``He's the president of the United States; he does have to lead and he will,'' said Senator Kent
Conrad, Democrat of North Dakota. ``But he's got to pick his spots.''

As lawmakers struggle to work out their differences, Mr.~Obama is courting the governors, an effort
that one White House official, speaking anonymously to discuss strategy, said began with last week's
meeting. In interviews, two governors -- a Democrat, Christine Gregoire of Washington, and a
Republican, Michael Rounds of South Dakota -- both said Mr.~Obama asked them to talk to members of
Congress about their own innovative approaches to health reform. Both said he urged them to be
careful about their language.

``I think he said what we have to do is not call it rationing, because clearly there is from H.M.O.
days a concern about rationing,'' Mr.~Rounds said, adding that he sensed Mr.~Obama was hoping to
have ``a bipartisan group of governors working directly with lawmakers to perhaps break a
stalemate.''

Ms.~Gregoire said the president reminded the governors that ``Congress has a bad taste in its mouth
from previous experience with managed care,'' and suggested they avoid the term. Instead, he spoke
of ``evidence-based care,'' the practice of using research to guide medical decisions. She said the
president told them, ``I need you to stick with me.''

Governors are deeply concerned about the rising price of Medicaid, the government insurance program
for the poor, which makes them natural allies of the president, who has made driving down health
costs a centerpiece of his effort.

But some of the proposals under consideration in Washington would expand Medicaid, a prospect that
the governors find worrisome. They also fear that any legislation passed by Congress would undermine
their own efforts at health reform, and said they used last week's meeting to tell Mr.~Obama so.

In reaching out to the governors, Mr.~Obama is reprising a strategy that worked for him during the
debate over his economic stimulus package, when he found far more support among Republican governors
than among Republican lawmakers.

Whether Mr.~Obama can have a more bipartisan outcome with health care remains unclear. He has
invested so much political capital in a health care bill that not to have legislation would be
politically disastrous for him. If that means passing a bill without Republican support, some
Democrats say, Mr.~Obama will do it.

``His instinct is to try, if he can, to find an honorable compromise with Republicans,'' Mr.~Podesta
said. ``But ultimately what he cares about at the end of the day is sitting there, pen in hand,
signing a bill that's a good bill, and that he believes in. When all is said and done, that's what I
think they will drive toward.''

\section{A 3,000-Mile Triumph, Spurrexd On by Diabetes}

Last week, a team of eight cyclists completed the coast-to-coast bike marathon called the Race
Across America in record time. It was quite an achievement under any circumstances, but what made it
extraordinary was something all eight of them had in common: Type 1 diabetes.

Type 1, sometimes called juvenile diabetes, poses special challenges for athletes. A person with
Type 1 can't produce insulin and must take regular injections to control blood sugar. But exercise
can also lead to precipitous, even deadly, drops in blood sugar. (Type 2 diabetes, by far the more
common form of the disease, typically develops later in life, often linked to poor eating habits and
weight gain; exercise is often prescribed as a way to keep blood sugar low.)

The accomplishments of the cyclists, who have a corporate sponsor and ride as Team Type 1, have
become a source of inspiration for the estimated three million Americans with Type 1 diabetes, and
especially for worried parents confronting a diagnosis of the disease in their children.

But the victory also offers lessons for the rest of us, underlining the benefits of daily vigilance
when it comes to health. Because people with Type 1 produce no insulin, they cannot survive without
injecting it before each meal, and they must wear a monitor or test their blood several times a day
to check their glucose levels. Meals, snacks, exercise and medication are carefully balanced. This
meticulous regimen is necessary to prevent diabetes complications, which can include kidney failure,
blindness and death. But closely controlling blood sugar can also result in an enviable level of
weight management and overall health.

``We're showing people that diabetes is our strength, and because of it we can do some pretty
amazing things,'' said Phil Southerland, 27, a cyclist and runner from Atlanta who founded the team
with another cyclist with Type 1, Joe Eldridge. (Both men rode in the race in 2006 and 2007, but did
not compete this year.)

``I think the rest of the world can look at the team,'' Mr.~Southerland continued, ``and say: 'These
guys just won a bike race, and they did it with diabetes. What can I do with my life to live a
healthier, better life?' ''

The achievements of the Type 1 athletes come at a time of growing concern about changing patterns of
the disease. While Type 2 diabetes is associated with an unhealthy lifestyle, scientists do not yet
know what causes Type 1, although autoimmune, genetic and environmental factors appear to play a
role.

With obesity and poor health habits rising among adults and children alike, it is hardly surprising
that Type 2 diabetes has grown so prevalent. But now European health officials are reporting an
inexplicable rise in Type 1 as well. Last month, the medical journal Lancet reported that the
incidence was rising about 4 percent a year among European children, particularly those under 5. At
that rate, the number of Type 1 cases will increase 70 percent in the next decade. And the disease
also appears to be gaining in the United States.

The rapid rise suggests environmental influences, and researchers are looking at possible factors
that include Caesarean deliveries, viral infection and early-life nutrition, including Vitamin D
deficiency.

Type 1 diabetes has many faces. Perhaps the most prominent spokeswoman is Mary Tyler Moore, now 72;
President Obama's nominee for a vacancy on the Supreme Court, Judge Sonia Sotomayor, 55, was given a
Type 1 diagnosis when she was 8.

And Jay Cutler, the star quarterback who was recently traded to the Chicago Bears, learned he had
Type 1 at 24, after a rapid 35-pound weight loss and severe fatigue originally attributed to stress.
A team physical finally led to a diagnosis.

``When I found out I think I was more relieved than anything else, just to know it could be
managed,'' Cutler said in an interview. ``You feel sorry for yourself for a little bit, but you come
to terms with it.''

Now when Cutler heads to the sidelines during practices or a game, he is met by a trainer who checks
his blood sugar to make sure it hasn't dropped to a risky level.

``The trainer is there with my meter, they prick it, do the whole thing, tell me my number,'' he
said. ``If we're good, we keep going. If I'm getting a little low I take some Gatorade.''

During the Race Across America, the Team Type 1 cyclists wore glucose monitors and traveled with a
doctor, eating or drinking when blood sugar levels begin to drop. While all the athletes must take
insulin regularly to prevent high blood sugar, the intense exercise causes their medication needs to
drop 60 to 75 percent during the first days of the race. As the body adjusts in the later days of
the race, the cyclists must generally increase their insulin injections.

To complete the race, the cyclists divided into two teams of four. The first four riders took turns
pedaling at full sprint for 10 to 15 minutes at a time, meaning each individual rider took only a
short break before having to ride again. After about 150 miles of tag-team racing, the exhausted
cyclists met up with the second set of four riders, who took over, giving the first riders time to
eat and rest before they started again.

Despite mechanical problems, the team, which started in Oceanside, Calif., rolled into Annapolis,
Md., a distance of 3,021 miles, in five days, nine hours and five minutes. Their average speed was
23.41 miles per hour -- 0.17 better than the winner last year, a Norwegian cycling team made up of professionals.

\section{Blink Twice if You Like Me}

LINCOLN, Mass. -- Sara Lewis is fluent in firefly. On this night she walks through a farm field in
eastern Massachusetts, watching the first fireflies of the evening rise into the air and begin to
blink on and off. Dr.~Lewis, an evolutionary ecologist at Tufts University, points out six species
in this meadow, each with its own pattern of flashes.

Along one edge of the meadow are Photinus greeni, with double pulses separated by three seconds of
darkness. Near a stream are Photinus ignitus, with a five-second delay between single pulses. And
near a forest are Pyractomena angulata, which make Dr.~Lewis's favorite flash pattern. ``It's like a
flickering orange rain,'' she said.

The fireflies flashing in the air are all males. Down in the grass, Dr.~Lewis points out, females
are sitting and observing. They look for flash patterns of males of their own species, and sometimes
they respond with a single flash of their own, always at a precise interval after the male's.
Dr.~Lewis takes out a penlight and clicks it twice, in perfect Photinus greeni. A female Photinus
greeni flashes back.

``Most people don't realize there's this call and response going on,'' Dr.~Lewis said. ``But it's
very, very easy to talk to fireflies.''

For Dr.~Lewis, this meadow is the stage for an invertebrate melodrama, full of passion and yearning,
of courtship duets and competitions for affection, of cruel deception and gruesome death. For the
past 16 years, Dr.~Lewis has been coming to this field to decipher the evolutionary forces at play
in this production, as fireflies have struggled to survive and spread their genes to the next
generation.

It was on a night much like this one in 1980 when Dr.~Lewis first came under the spell of fireflies.
She was in graduate school at Duke University, studying coral reef fish. Waiting for a grant to come
through for a trip to Belize, she did not have much else to do but sit in her backyard in North
Carolina.

``Every evening there was this incredible display of fireflies,'' Dr.~Lewis said. She eventually
started to explore the yard, inspecting the males and females. ``What really struck me was that in
this one-acre area there were hundreds of males and I could only find two or three females,'' she
said. ``I thought, 'Man, this is so intense.' ''

When a lot of males are competing for the chance to mate with females, a species experiences a
special kind of evolution. If males have certain traits that make them attractive to females, they
will mate more than other males. And that preference may mean that those attractive males can pass
down their traits to the next generation. Over thousands of generations, the entire species may be
transformed.

Charles Darwin described this process, which he called sexual selection, in 1871, using male
displays of antlers and feathers as examples. He did not mention fireflies. In fact, fireflies
remained fairly mysterious for another century. It was not until the 1960s that James Lloyd, a
University of Florida biologist, deciphered the call and response of several species of North
American firefly.

Dr.~Lewis, realizing that other firefly mysteries remained to be solved, switched to fireflies from
fish in 1984, when she became a postdoctoral researcher at Harvard. She taught herself Dr.~Lloyd's
firefly code and then began to investigate firefly mating habits. North American fireflies spend two
years underground as larvae, then spend the final two weeks of their lives as adults, flashing,
mating and laying eggs. When Dr.~Lewis started studying fireflies, scientists could not say whether
the females mated once and then laid all their eggs, or mated with many males. ``Nobody knew what
happened after the lights went out,'' Dr.~Lewis said.

She searched for mating fireflies in the evening, marked their locations with surveyor's flags and
then revisited them every half-hour through the night. They were still mating at dawn.

``It was cool to watch the sun rise and see the couples breaking up and the females crawling down
the grass to lay their eggs,'' Dr.~Lewis said.

Many Americans are familiar with the kinds of fireflies Dr.~Lewis studies, but they represent only a
tiny fraction of the 2,000 species worldwide. And there is enormous variation in these insects.
``There are some species that produce flashes when they're adults, and there are some that simply
glow as adults,'' Dr.~Lewis said. ``Then there are a whole bunch of species where the adults don't
produce any light at all.''

In recent years scientists have analyzed the DNA of fireflies to figure out how their light has
evolved. The common ancestor of today's fireflies probably produced light only when they were
larvae. All firefly larvae still glow today, as a warning to would-be predators. The larvae produce
bitter chemicals that make them an unpleasant meal.

As adults, the earliest fireflies probably communicated with chemical signals, the way some firefly
species do today. Only much later did some firefly species gain through evolution the ability to
make light as adults. Instead of a warning, the light became a mating call. (An enzyme in the
firefly's tail drives a chemical reaction that makes light.)

The more Dr.~Lewis watched firefly courtship, the clearer it became that the females were carefully
choosing mates. They start dialogues with up to 10 males in a single evening and can keep several
conversations going at once. But a female mates with only one male, typically the one she has
responded to the most.

Dr.~Lewis wondered if the female fireflies were picking their mates based on variations in the
flashes of the males. To test that possibility, she took female fireflies to her lab, where she has
computer-controlled light systems that can mimic firefly flashes. ``You can play back specific
signals to females and see what they respond to,'' Dr.~Lewis said.

The female fireflies turned out to be remarkably picky. In many cases, a male flash got no response
at all. In some species, females preferred faster pulse rates. In others, the females preferred
males that made long-lasting pulses.

If females preferred some flashes over others, Dr.~Lewis wondered why those preferences had evolved
in the first place. One possible explanation was that the signals gave female fireflies a valuable
clue about the males. Somehow, mating with males with certain flash patterns allowed females to
produce more offspring, which would inherit their preference.

It is possible that females use flashes to figure out which males can offer the best gifts. In many
invertebrate species, the males provide females with food to help nourish their eggs. Dr.~Lewis and
her colleagues discovered that fireflies also made these so-called nuptial gifts -- packages of
protein they inject with their sperm.

Dr.~Lewis is not sure why she and her colleagues were the first to find them. The gifts form coils
that can take up a lot of space in a male firefly's abdomen. ``They're incredibly beautiful,'' she
said.

Receiving nuptial gifts, Dr.~Lewis and her colleagues have shown, can make a huge difference in the
reproductive success of a female firefly. ``It just about doubles the number of eggs a female can
lay in her lifetime,'' she said. One reason the effect is so big is that fireflies do not eat during
their two-week adulthood. A slowly starving female can use a nuptial gift to build more eggs.

In at least some species, females may use flashes to pick out males with the biggest gifts.
Dr.~Lewis has tested this hypothesis in two species; in one, males with conspicuous flashes have
bigger gifts. In another species, she found no link.

``In some cases they could be honest signals, and in some cases they could be deceptive signals,''
Dr.~Lewis said.

Deception may, in fact, evolve very easily among fireflies. It turns out that a male firefly does
not need to burn many extra calories to make flashes. ``It takes some energy, but it's tiny. It's
less costly for a male than flying around,'' Dr.~Lewis said.

If making light is so cheap for males, it seems odd that they have not all evolved to be more
attractive to females. ``What is it that keeps their flashes from getting longer and longer or
faster and faster?'' Dr.~Lewis asked.

Scanning the meadow, she grabbed her insect net and ran after a fast-flying firefly with a triple
flash. She caught an animal that may offer the answer to her question. Dr.~Lewis dropped the insect
into a tube and switched on a headlamp to show her catch. Called Photuris, it is a firefly that eats
other fireflies.

``They are really nasty predators,'' Dr.~Lewis said. Photuris fireflies sometimes stage aerial
assaults, picking out other species by their flashes and swooping down to attack. In other cases,
they sit on a blade of grass, responding to male fireflies with deceptive flashes. When the males
approach, Photuris grabs them.

``They pounce, they bite, they suck blood -- all the gory stuff,'' Dr.~Lewis said. She has found
that each Photuris can eat several fireflies in a night. Photuris kills other fireflies only to
retrieve bad-tasting chemicals from their bodies, which it uses to protect itself from predators.

To study how Photuris predation affects its firefly prey, Dr.~Lewis and her colleagues built sticky
traps equipped with lights that mimicked courtship signals of Photuris's victims. The scientists
found that Photuris was more likely to attack when flash rates were faster. In other words,
conspicuous flashes -- the ones females prefer -- also make males more likely to be killed.

``At least where Photuris predators are around,'' Dr.~Lewis said, ``there's going to be a strong
selection for less conspicuous flashes.''

\section{The Roberts Court, Tipped by Kennedy}

WASHINGTON -- Chief Justice John G.~Roberts Jr.~emerged as a canny strategist at the Supreme Court
this term, laying the groundwork for bold changes that could take the court to the right even as the
recent elections moved the nation to the left.

The court took mainly incremental steps in major cases concerning voting rights, employment
discrimination, criminal procedure and campaign finance. But the chief justice's fingerprints were
on all of them, and he left clues that the court is only one decision away from fundamental change
in many areas of the law.

Whether he will succeed depends on Justice Anthony M.~Kennedy, the court's swing vote. And there is
reason to think that the chief justice has found a reliable ally when it counts.

``In the important cases, Kennedy ends up on the right,'' said Thomas C.~Goldstein, a student of the
court and the founder of Scotusblog, which has compiled comprehensive statistics on the current
term. The two justices agreed 86 percent of the time.

If Judge Sonia Sotomayor is confirmed by the Senate, she will succeed Justice David H.~Souter, a
liberal who spent almost two decades on the court. Her record on the federal appeals court in New
York suggests that her views are largely in sync with those of Justice Souter, though there is some
evidence that she will turn out to be more conservative in criminal cases.

The arrival of a neophyte justice coupled with Chief Justice Roberts's increasing mastery of the
judicial machinery foreshadow a widening gap between the Democratic-led political branches and the
Supreme Court. Indeed, the court appears poised to move to the right in the Obama era.

Chief Justice Roberts has certainly been planting seeds in this term's decisions. If his reasoning
takes root in future cases, the law will move in a conservative direction on questions as varied as
what kinds of evidence may be used against criminal defendants and the role the government may play
in combating race discrimination.

The two newest justices, Chief Justice Roberts and Justice Samuel A.~Alito Jr., both appointed by
President George W.~Bush, agreed 92 percent of time, the highest rate for any pair of justices. But
Justice Alito often wrote concurring opinions to underscore or try to extend conservative rulings,
especially in criminal cases. He may well now be the court's most conservative member.

``Alito is staking out some room to the right of the chief justice,'' said Pamela Harris, the
executive director of the Supreme Court Institute at Georgetown University Law Center, ``and you
would have thought there is no such room.''

Labels like ``conservative'' and ``liberal'' are, of course, imperfect. Political scientists often
say judges are liberal to the extent they tend to vote in favor of, say, criminal defendants,
environmental groups, people suing over injuries and plaintiffs claiming discrimination.

Decisions protecting individual rights may be said to be liberal, too, but many political liberals
would reject that characterization where the First Amendment rights of rich campaign contributors or
the Second Amendment rights of gun owners are at issue.

At the Supreme Court, though, voting alignments are so predictable that ``liberal'' and
``conservative'' are as much shorthand as principle. They refer to two groups, of four justices
each, who generally vote the same way.

The court was remarkably polarized in the 74 signed decisions it issued this term, dividing 5-to-4
or 6-to-3 in almost half of them, up from roughly a third in the three previous years. The court
reversed lower courts about three-quarters of the time, up from two-thirds in the last term.

Justice Kennedy was in the majority 92 percent of the time and in all but 5 of the 23 decisions in
which the justices split 5-to-4. Those decisions were, moreover, often divided in the expected way:
in 16, all four members of the court's liberal wing were on one side and all four of its
conservatives were on the other.

And in between them was Justice Kennedy, the most powerful jurist in America. He joined the liberals
5 times and the conservatives 11. That was a significant shift to the right: in the previous term,
Justice Kennedy voted four times each with the liberals and the conservatives in cases divided along
the traditional ideological fault line.

Justice Kennedy swung right in the cases that really mattered. The only major case in which he
joined the court's liberal wing -- Justices Souter, John Paul Stevens, Ruth Bader Ginsburg and
Stephen G.~Breyer -- was Caperton v. A.T. Massey Coal Company. The decision required the chief
justice of the West Virginia Supreme Court to disqualify himself from a case involving a coal
executive who had spent \$3 million to elect him.

If there were surprises, they came from Justices Antonin Scalia and Clarence Thomas.

``For all the talk about Scalia and Thomas being the most conservative justices on the court, they
are the justices most likely in play,'' said Jeffrey L.~Fisher, a law professor at Stanford who has
argued several important criminal cases before the court.

Justices Scalia and Thomas are apt to follow what they understand to be the original meaning of the
Constitution, even when the consequences might not align with their policy preferences. In
Melendez-Diaz v. Massachusetts, for instance, Justices Scalia and Thomas joined three members of the
court's liberal wing to say that the Constitution's confrontation clause requires crime laboratory
analysts to appear at trial rather than submit written reports.

The current chief justice clerked for Chief Justice William H.~Rehnquist, a famous strategist, and
he seems to have learned some tactics from his old boss.

The most important case of the term, for instance, seemed likely to be Northwest Austin Municipal
Utility District v. Holder, which concerned the constitutionality of a major provision of the Voting
Rights Act of 1965. Judging from the questions at argument, Chief Justice Roberts appeared prepared
to strike it down.

The law, reauthorized by Congress in 2006, used old data to decide which state and local governments
had to obtain federal permission before making even minor changes to their voting procedures. Some
officials, mostly in the South, bristled at what they said was the stigma, burden and federal
intrusion that came with being covered by the law.

Instead of addressing that broad question, Chief Justice Roberts wrote a narrow decision for an
eight-justice majority that allowed the Texas municipal water district that had brought the
challenge to try to escape from coverage under the law while leaving the law itself intact.

In the process, though, he wrote that much has changed since the civil rights era and suggested that
the law would not survive constitutional scrutiny. He all but invited a further challenge, and many
voting rights specialists say that Congress must act soon to change the law if it is to survive
another encounter with the Roberts court.

What accounted for the incrementalism? A likely explanation is that the chief justice did not yet
have Justice Kennedy's unqualified support and was biding his time until he did.

Something similar seemed to be going on in Herring v. United States, which cut back on but did not
eliminate the exclusionary rule. The rule requires the suppression of some evidence obtained by
police misconduct.

The majority opinion in Herring, by Chief Justice Roberts, said that an unlawful arrest based on
incorrect information in a police database did not require suppression of drugs and a gun. The
police error, the chief justice wrote, was ``isolated negligence attenuated from the arrest.''

The balance of the opinion suggested that Chief Justice Roberts might be willing to excuse other
kinds of police misconduct as well, and specialists in criminal law said they suspected the limiting
principle was added to placate Justice Kennedy, who in a 2006 concurrence expressed support for the
exclusionary rule even as he joined a majority opinion that harshly criticized it.

Even the court's failure to decide a case could seem portentous.

In scheduling Citizens United v. Federal Election Commission for re-argument in September, Chief
Justice Roberts appeared to be setting the stage for an overhaul of the law governing campaign
spending by corporations.

The case involves a polemical documentary about Hillary Rodham Clinton that the F.E.C. said was an
``electioneering communication'' that could not be broadcast during the Democratic primaries this
year because a corporation had financed it.

The Supreme Court could have ruled in favor of the group that made the film on a variety of narrow
grounds. Instead, it asked for fresh briefing on the validity of laws limiting corporate campaign
spending.

In two of the term's most important cases, Justice Kennedy was the on-off switch in determining the
meaning of the Constitution's due process clause.

In District Attorney's Office v. Osborne, Justice Kennedy was not prepared to find that due process
requires allowing inmates to have access to DNA evidence that might prove their innocence, and so
there is no such right in the Constitution.

In the Caperton decision, he said litigants' due process rights required a state supreme court
justice to step aside from a case involving an important supporter, and now that has become a
constitutional principle.

The court's liberals would have interpreted the due process clause to require both kinds of rights,
the court's conservatives neither. The Constitution, it turns out, means what Justice Kennedy says it means.

\section{In Skittish Hollywood, Stars Can't Save 'Moneyball'}

LOS ANGELES -- In a film production office here, at least a couple of employees were still hanging
around on Monday, hoping in vain to score with their troubled baseball movie ``Moneyball.''

But they had swung, and missed.

Just days before shooting was to begin, Sony Pictures pulled the plug on ``Moneyball,'' a major film
project starring Brad Pitt and being directed by Steven Soderbergh. The last-minute demise of a
prestige picture is a rare spectacle in Hollywood -- one that is painful, expensive and damaging to
all involved. This one is estimated to have cost Sony \$10 million in script development and costs
like scouting locations.

But such disasters may become more common as an increasingly nervous film business comes to terms
with a sharp decline in home video revenue and with the diminishing leverage of even the most
popular A-list stars, like Mr.~Pitt.

``They're much more careful about doing a movie just because a star wants to do it,'' said Eric
Weissmann, a longtime entertainment lawyer who recalled the days when Warner Brothers made a film,
``An Enemy of the People,'' based on an Ibsen play, largely because Steve McQueen wanted to do it.

Representatives of Sony, Mr.~Pitt and Mr.~Soderbergh all declined to discuss ``Moneyball.'' But
accounts from more than a dozen people involved with the film, who spoke on the condition of
anonymity to avoid damaging professional relationships, described a process in which the heady rush
toward production was halted by a studio suddenly confronted by plans for something artier and more
complex than bargained for.

The swift mothballing of ``Moneyball'' may also increase doubts that Hollywood can still deliver
tricky but appealing pictures like ``Michael Clayton,'' ``Good Night, and Good Luck'' or ``The
Curious Case of Benjamin Button.'' Studio divisions that specialized in sophisticated films, like
Warner Independent Pictures and Paramount Vantage, have recently been closed or diminished.

A central player in the behind-the-scenes drama over ``Moneyball'' was Amy Pascal, Sony Pictures'
co-chairman and an executive known for taking a strong hand in the development of scripts.
Ms.~Pascal and her team became involved with ``Moneyball'' about six years ago, when a relatively
untested producer, Rachael Horovitz, brought the idea to Sony with a screenwriter, Stan Chervin,
after virtually every other buyer in Hollywood had passed.

``Moneyball,'' which is based on a 2003 nonfiction book by Michael Lewis, tells the story of Billy
Beane, the Oakland Athletics general manager who figured out how to build a winning team on the
cheap with players undervalued by the conventional measures of success in baseball.

With a budget estimated at about \$57 million, it was not hugely expensive but not a small indie
project, either. The film was of a sophisticated type that needed the cachet of a Soderbergh, the
star power of a Pitt and perhaps Academy Award potential to overcome its somewhat cerebral quality
and the difficulty of attracting foreign viewers for a movie about America's pastime.

Two writers who had worked with Sony on ``Ali,'' Stephen Rivele and Christopher Wilkinson, did
script work after Mr.~Chervin's version, which tried to warm up the story by focusing on the
relationship between Mr.~Beane and his daughter. Mr.~Chervin returned to the project to work with
the director David Frankel, who chose to do ``Marley \& Me'' instead.

Mr.~Pitt, a fan of the book, had become interested, putting the film on a fast track at Sony. The
studio hired Steven Zaillian, a reliable screenwriter who had won an Oscar for ``Schindler's List,''
to do another rewrite, even as it agreed to bring on Mr.~Soderbergh as the director.

Although he has scored big with studio projects like the ``Ocean's Eleven'' series with Mr.~Pitt,
Mr.~Soderbergh remains one of Hollywood's most self-consciously distinctive directors. He serves as
his own cinematographer, often contributes to scripts and has worked lately on a series of
challenging, low-budget films like his two-part ``Che,'' a Spanish-language movie that made its
debut both in a small number of theaters and on pay-per-view cable TV.

Two weeks ago, a mismatch in personal style and expectations collided. Mr.~Soderbergh, a week before
filming was to begin, delivered his own revision of the script.

One reason was to win the approval of Major League Baseball, which was not happy with some factual
liberties in Mr.~Zaillian's version. Such approval is crucial in a baseball film that intends to use
protected trademarks.

``Typically, on a film like this, we look at it for historical accuracy,'' said Matthew Bourne, a
vice president of Major League Baseball for public relations. ``We've been in touch with Soderbergh
and Sony, and they've been receptive to our requests.''

What baseball saw as accurate, Sony executives saw as being too much a documentary. Mr.~Soderbergh,
for instance, planned to film interviews with some of the people who were connected to the film's
story.

The executives, who had just seen disappointing results from ``The Taking of Pelham 1 2 3'' and
``Year One,'' rebelled. Ms.~Pascal and Matt Tolmach, co-president of Sony's film operation,
personally told Mr.~Soderbergh of their dissatisfaction.

The situation was particularly ticklish, given Ms.~Pascal's close professional relationship with
Bryan Lourd, the Creative Artists partner who serves as one of Mr.~Pitt's agents. In a highly
unusual arrangement, when the studio decided to pull the plug, it allowed representatives for
Mr.~Soderbergh, Mr.~Pitt and the producers a weekend to offer the film to Paramount and Warner. Both
studios, however, immediately passed.

Through last week, the ``Moneyball'' team looked for a compromise that might restart the film. But
Fox, which also got a look, decided to pass as well.

By this week, the movie, at least in its current configuration, was dead. Mr.~Pitt's representatives
had an eye out for his next picture. Mr.~Soderbergh's were looking for ways to assure that his
valuable, if somewhat eccentric, career would not be harmed by the debacle.

As of Tuesday, ``Moneyball'' was back in development, with Sony executives still hoping at some
point to work with Mr.~Pitt and to find a director more interested in Mr.~Zaillian's version.

``There's a movie in there,'' Mr.~Wilkinson said on Monday. ``But it's a very unusual movie.''

\section{Hong Kong's Pro-Democracy March Draws Thousands}

HONG KONG -- Thousands of people joined a pro-democracy march here on Wednesday, although the
turnout fell short of a candlelight vigil held nearly four weeks ago to commemorate the 20th
anniversary of the Tiananmen Square crackdown in Beijing.

An enormous crowd for the annual June 4 candlelight vigil, the largest since 1990, had raised the
hopes of Hong Kong democracy advocates that the same enthusiasm might carry over to their movement.
The movement has been struggling after several small successes from 2003 to 2005, including winning
support for blocking the government's planned introduction of stringent internal security
legislation.

The immediacy of democracy demands here has faded somewhat as Beijing officials have ruled out
direct elections for the chief executive until 2017 and the legislature until 2020.

The march on Wednesday, on the 12th anniversary of Hong Kong's return to Chinese rule after 156
years of British control, nonetheless drew a large crowd.

Many marchers said they were dissatisfied with government policies to deal with the economy.
Unemployment in Hong Kong rose sharply over the winter and leveled off this spring at 5.3 percent --
a little over half the rate in the United States, but a shock for a territory where the rate was 3.2
percent last summer.

But the largest single issue seemed to be the limits on democracy in Hong Kong. ``The majority comes
here for democracy, but there are other grievances against government policy,'' said Sin Chung Kai,
vice chairman of the Democratic Party.

When Britain returned Hong Kong to Chinese rule in 1997, the Chinese government initially held out
the possibility of full democracy after 2007, including the concept in the Basic Law, Hong Kong's
miniconstitution, but stopped short of an unequivocal promise of how and when to achieve universal
suffrage.

A committee of 800 people, most with connections to Beijing, chooses the chief executive here, who
must then be appointed by leaders in Beijing before taking office. Half the legislature is chosen by
the public and half by a variety of interest groups, including banks, chambers of commerce, trade
unions and lawyers.

The police estimated that 26,000 people had assembled in Victoria Park on Hong Kong Island as the
march began. The organizers had said that they expected more to join the march along the way, and
they estimated that 76,000 people took part.

The police had estimated the crowd at the June 4 Tiananmen vigil, at the same location in Victoria
Park, at 62,800, while organizers put it at 150,000.

The vigil did have some carryover effect on Wednesday's march. Jupiter Chan, a 24-year-old graduate
student, said that the vigil prompted him to come to the annual democracy march this year for the
first time since 2003.

``I was touched by the Fourth of June ceremony, and I felt that if I didn't come this year, I would
regret it later,'' he said.

\section{Rural South Koreans' Global Links Grow, Nourished by a Satellite Crop}

YEONGJU, South Korea -- Lee Si-kap, a shy farmer living in this central South Korean town, holds a
record: He owns more satellite dishes than any other South Korean -- 85 of them, receiving 1,500
satellite television channels from more than 100 countries, some as far away as South Africa and
Canada.

To passers-by, Mr.~Lee's home stands like an exclamation mark in the otherwise nondescript
countryside dotted with apple orchards and ginseng fields. Satellite dishes cover his roof like
giant steel mushrooms.

They spread into his front yard and blossom in a field behind his house, some as large as 16 feet in
diameter.

Once dismissed as a local eccentric, Mr.~Lee has more recently emerged as something of a hero of
modest fame, featured on national television as ``antenna man.'' Since late last year, he and
thousands of fellow satellite enthusiasts -- including the husbands of foreign brides and a few
dedicated souls searching for signals from extraterrestrial life forms -- have started a campaign to
install free satellite dishes for poor foreign brides living in rural South Korea, so they can
receive broadcasts from their home countries.

``Thanks to Mr.~Lee, I now miss my country, my mother and father less than I used to,'' said Bui Thi
Huang, a 22-year-old bride from Haiphong, Vietnam, who now lives here in Yeongju, about 100 miles
southeast of Seoul.

In recent years, the South Korean countryside has had an influx of brides from poorer countries like
Vietnam, China and the Philippines. Like Ms.~Bui, they marry South Korean farmers who have
difficulty finding a spouse because so many young Korean women have rejected rural life and migrated
to cities.

In towns like Yeongju, these young foreign brides have become a bedrock of the local economy. They
work alongside their husbands in the fields and have brought back a sound that was fast becoming a
distant memory among the aging farm population here: crying babies.

In South Korea, which had once prided itself on being a homogeneous society, 4 out of 10 women who
married in rural communities last year were foreign born. In Yeongju alone, the number of foreign
wives increased by 28 percent in the past year and a half, to 250, half of them from Vietnam.

``These women have a hard time fitting in. The local governments, and the husbands, often focus only
on making them 'Korean,' teaching them the Korean language and computer skills,'' said Mr.~Lee, 39,
who has never married. ``They don't quite understand how isolated these women feel.''

When Mr.~Lee, who lives with his 80-year-old mother and 97-year-old grandfather, is not toying with
his satellite equipment, he tends his pepper and sesame fields or makes the rounds of nearby
villages to see if the foreign brides are having any problems with their television reception.

Mr.~Lee and his friends still encounter objections from husbands who are determined to shield their
foreign brides from any reminders of their native lands, for fear these might only magnify their
homesickness. But they are encouraged that many families have reported that watching satellite
broadcasts from home actually helps the women to overcome their loneliness and better adjust to life
here.

Mr.~Lee says his sympathy for foreign brides stems in part from his own experience of feeling cut
off from society.

He felt deeply hurt when his father abandoned him and his mother when he was a small boy, and,
lacking self-confidence, had trouble making friends in his neighborhood and at school. He rarely
ventured outside his village, and said he still feared making phone calls.

What saved him, he said, was ``music -- and satellite television.''

``Music was my only friend,'' said Mr.~Lee, whose dream is to meet his idol, the American heavy
metal rock musician Ronnie James Dio. ``And because I couldn't get much rock music on Korean
television, I turned to satellite television.''

Satellite television introduced him to the wider world -- to Japanese baseball, life on Pacific
islands, Russian folk music and religions in India and Nepal.

He installed his first satellite dish in 1992, when he was 23 and had already returned to farming
after receiving a vocational college degree in electronics. Collecting secondhand satellite dishes
has since become a hobby that verges on an obsession. When most farmers here look to the sky, they
read clouds for weather. When Mr.~Lee looks skyward, he says, he imagines satellites in earth orbit.
To him, the air is filled with broadcast signals, ``like seeds from thistles.''

Farmers here at first did not know what to make of their bachelor neighbor, who listened to heavy
metal music, often belting out the lyrics in English, sometimes in Japanese. They would see him on
the roof under the blazing sun of summer or under the starry winter sky, fiddling for hours with his
satellite equipment.

Although he does not understand most languages on the broadcasts he receives, Mr.~Lee said: ``It
gets addictive. The more dishes you have, the more channels you can get.

``Nothing compares with the joy of catching a new broadcast channel from a faraway country,'' he
said. ``It's like pulling in a big fish. It's the excitement of discovering something from outside
the boundaries of your usual world.''

\section{From a Beijing Suburb, Vibrant Strings}

DONGGAOCUN, China -- Perhaps the only thing more aurally challenging than a roomful of novice
violinists screeching their way through ``Mary Had a Little Lamb'' is a roomful of novice violinists
screeching along on out-of-tune instruments.

``Stop,'' Chen Yiming shouted to her students, an enthusiastic bunch, ages 8 to 47. ``Can we please
pay attention to our instruments and make sure they are tuned correctly?''

After a short break for adjustments, the cacophony resumed.

Violin fever has hit this drab rural township with hundreds of residents, young and old, picking up
the bow as Donggaocun tries to position itself as the string instrument capital of China.

Once known primarily for its abundant peach harvest, the town, about an hour's drive from downtown
Beijing, has become one of the world's most prodigious manufacturers of inexpensive cellos, violas,
violins and double basses. Last year the town's 9 factories and 150 small workshops made 250,000
instruments, most of them ending up in the hands of students in the United States, Britain and
Germany.

The city fathers have taken to boasting that Donggaocun produces 30 percent of the world's string
instruments, although another town in southeast China, Xiqiao, makes a similar, if slightly more
credible, claim.

Feng Yuankai, of the China Musical Instrument Association, said that Xiqiao, in Jiangsu Province, is
still the top producer, but that Donggaocun, which opened its first factory in 1988, is catching up.
``Even with the economic slowdown their output is growing very quickly,'' he said, adding that
Donggaocun is responsible for less than 20 percent of China's violin market.

A disputed ranking, or the recent slump in orders, is not about to dent Donggaocun's ambition of
becoming China's violin mecca. Last month, officials began promoting the creation of a tourist
attraction they are calling China's first ``string instrument experience center.'' When it opens
this summer, Instrument City will include a museum of ``world-famous'' violins, a 500-seat concert
hall, a musical fountain and what the vice mayor describes as a ``folklore village.''

Then there was this spring's annual peach blossom festival, which featured a violin extravaganza
with 1,000 fiddlers taking to the stage. Even if about half of the players -- those too green for
public consumption -- were asked to refrain from actually playing, the concert was an auditory tour
de force.

Since 2006 the local education department has trained 40 teachers to become violinists so that every
school in Donggaocun and surrounding communities can offer music instruction. Violins not only have
become a driver of economic development, but also have elevated the town's sense of itself, said
Wang Junying, the department's chief propaganda officer. ``They have helped us become a more
cultured and elegant place,'' Ms.~Wang said.

The workers at Beijing Hongsheng Yun Violin Instruments Company largely agree. Most of the factory's
employees used to work in the same building, a paper plant that belched out noxious smoke. But in
their effort to clean Beijing's air for the Summer Olympics last August, the authorities closed the
plant and invested government money in violin production. Although 200 workers lost their job when
the paper plant closed, the two dozen people who were hired back by the violin workshop earn nearly
twice as much as they did before.

Among the workplace perks are free instruments for the children of employees. ``Violins have made us
richer, and they have raised our artistic awareness,'' said Zhao Gangcai, who assembles violins six
days a week and whose 13-year-old daughter recently started playing. ``Her classmates think it's
cool, and now they want to learn, too.''

Some of Donggaocun's most promising musicians end up in classes led by Ms.~Chen, 28, a woman of
inordinate patience who started playing the violin around the same time she began to talk. She
rejects the reluctant and those with stubby fingers, although she does take a handful of adults
whose zeal makes up for their lack of raw talent.

Few students match the eagerness of Song Wei, 47, a retired kindergarten teacher who gave up the
accordion for the violin after her husband and uncle began making string instruments. She practices
an hour or two each day, and several friends, also retirees, have also caught the violin bug,
although they dare not play together. ``The noise would be unbearable,'' she said with a laugh. ``My
only goal is to play a nice song when I hear it and to make myself happy.''

Ms.~Chen has loftier goals. She hopes to turn Donggaocun into a talent factory, although she
cautioned that prodigies usually emerge from musically inclined families, which, she suggested, were
in short supply in this largely rural township of 33,000. ``I want this place to produce world-class
musicians,'' she said.

In the meantime, she is working on her 1-year-old daughter. A few weeks ago, she gingerly placed a
violin in the girl's hands, and the results were encouraging. ``She knows how to hold it now,''
Ms.~Chen said. ``I'm taking that as a good sign.''

\section{Riots in Western China Amid Ethnic Tension}

BEIJING -- At least 1,000 rioters clashed with the police on Sunday in a regional capital in western
China after days of rising tensions between Muslim Uighurs and Han Chinese, according to witnesses
and photographs of the riot.

The rioting broke out Sunday afternoon in a large market area of Urumqi, the capital of the vast,
restive desert region of Xinjiang, and lasted for several hours before riot police officers and
paramilitary or military troops locked down the Uighur quarter of the city. The rioters threw stones
at the police and set vehicles on fire, sending plumes of smoke into the sky, while police officers
used firehoses and batons to beat back rioters and detain Uighurs who appeared to be leaders of the
protest, witnesses said.

At least three Han Chinese were killed in the rioting and 20 people were injured, according to
Xinhua, the official news agency. Dozens of Uighur men were led into nearby police stations with
their hands behind their backs and shirts pulled over their heads, one witness said. Early Monday,
the local government announced a curfew banning all traffic in the city until 8 p.m.

The riot was the largest ethnic clash in China since the Tibetan uprising of March 2008, and perhaps
the biggest protest in Xinjiang in years. Like the Tibetan unrest, it highlighted the deep-seated
frustrations felt by some ethnic minorities in western China over the policies of the Communist
Party.

Many Uighurs, a Turkic-speaking Muslim group, resent rule by the Han Chinese, and Chinese security
forces have tried to keep oil-rich Xinjiang under tight control since the 1990s, when cities there
were struck by waves of protests, riots and bombings. Last summer, attacks on security forces took
place in several cities in Xinjiang; the Chinese government blamed separatist groups.

Uighurs taking part in the riot on Sunday were furious over an ethnic brawl between Uighur and Han
workers that broke out on June 26 in a toy factory in Guangdong and that resulted in the deaths of
two Uighurs. A total of 116 people were injured. The police later arrested an ex-employee of the
factory who had ignited the brawl by starting a rumor that six Uighur men had raped two Han women at
the work site, Xinhua reported.

There was also a rumor going around on Sunday in Urumqi that a Han man had killed a Uighur earlier
that day in the city, said Adam Grode, an English teacher and former Fulbright scholar living in the
neighborhood where the rioting took place.

``This is just crazy,'' Mr.~Grode said in a telephone interview on Sunday night. ``There was a lot
of tear gas in the streets, and I almost couldn't get back to my apartment. There's a huge police
presence.''

Rumors of Uighurs attacking Han Chinese spread quickly through parts of Urumqi, adding to the panic.
A worker at Texas Restaurant, a few hundred yards from the site of the rioting, said her manager
urged the restaurant workers to stay inside. ``My boss went home in the evening and called us saying
he had heard that Uighurs were beating Han Chinese, so we'd better stay in the restaurant,'' said
the worker, a woman who gave only her surname, Wang.

Xinhua reported few details of the riot on Sunday night. It said that ``an unknown number of people
gathered Sunday afternoon'' in Urumqi, ``attacking passers-by and setting fire to vehicles.''

Uighurs are the largest ethnic group in Xinjiang but are a minority in Urumqi, where Han Chinese
make up more than 70 percent of the two million or so people. The Chinese government has encouraged
Han migration to the city and other parts of Xinjiang, fueling resentment among the Uighurs. Urumqi
is a deeply segregated city, with Han Chinese there rarely venturing into the Uighur quarter and
often warning visitors to stay away from the area.

The Uighur neighborhood is centered in a warren of narrow alleyways, food markets and a large
shopping area called the Grand Bazaar or the Erdaoqiao Market, where the rioting reached its peak on
Sunday.

Mr.~Grode, who lives in an apartment there, said he went outside when he first heard commotion
around 6 p.m. He saw hundreds of Uighurs in the streets; that quickly swelled to more than 1,000, he
said. When public buses stopped at the scene, Uighurs riding inside opened the windows and joined
the crowd in shouting slogans.

Police officers soon arrived. Around 7 p.m., protesters began hurling rocks, vegetables and other
goods from the market at the police, Mr.~Grode said. Traffic had ground to a halt, and some rioters
threw stones at bus windows.

An hour later, as the riot surged toward the center of the market, troops in green uniforms and full
riot gear showed up, as did armored vehicles. Chinese government officials often deploy the People's
Armed Police, a paramilitary force, to quell riots. The troops shot off tear gas canisters and might
have fired other projectiles too, Mr.~Grode said.

By midnight, he said, some of the armored vehicles had begun to leave, but bursts of gunfire could still be heard.

\section{Scores Killed in Clashes in Western China}

BEIJING -- The Chinese state news agency reported Monday that at least 140 people were killed and
816 injured when rioters clashed with the police in a regional capital in western China after days
of rising tensions between Muslim Uighurs and Han Chinese.

The casualty toll, if confirmed, would make this the deadliest outbreak of violence in China in many
years.

The rioting broke out Sunday afternoon in a large market area of Urumqi, the capital of the vast,
restive desert region of Xinjiang, and lasted for several hours before riot police officers and
paramilitary or military troops locked down the Uighur quarter of the city, according to witnesses
and photographs of the riot.

At least 1,000 rioters took to the streets, throwing stones at the police and setting vehicles on
fire. Plumes of smoke billowed into the sky, while police officers used fire hoses and batons to
beat back rioters and detain Uighurs who appeared to be leading the protest, witnesses said.

The casualty numbers appeared to be murky and shifting on Monday. A one-line report by Xinhua, the
state news agency, giving the estimate of 129 dead and 816 injured attributed the numbers to the
regional police department, but did not quote officials by name and did not have any details.
Earlier, Xinhua had reported that three civilians and one police officer had been killed.

One regional official reached by telephone put the death toll at 105 and said at least 800 people
had been injured. One American who watched the rioting at its height said he did not see people
being killed or corpses in the streets, though he said he did see Uighurs shoving or kicking a few
Han Chinese. Images of the rioting on state television showed some bloody people lying in the
streets and cars burning.

Dozens of Uighur men were led into police stations on Sunday evening with their hands behind their
backs and shirts pulled over their heads, one witness said. Early Monday, the local government
announced a curfew banning all traffic in the city until 8 p.m.

The riot was the largest ethnic clash in China since the Tibetan uprising of March 2008, and perhaps
the biggest protest in Xinjiang in years. Like the Tibetan unrest, it highlighted the deep-seated
frustrations felt by some ethnic minorities in western China over the policies of the Communist
Party, and how that can quickly turn into ethnic violence. Last year, in Lhasa, the Tibetan capital,
at least 19 people were killed, most of them Han civilians, according to government statistics.

Many Uighurs, a Turkic-speaking Muslim group, resent rule by the Han Chinese, and Chinese security
forces have tried to keep oil-rich Xinjiang under tight control since the 1990s, when cities there
were struck by waves of protests, riots and bombings. Last summer, attacks on security forces took
place in several cities in Xinjiang; the Chinese government blamed separatist groups.

Early Monday, Chinese officials said the latest riots were started by Rebiya Kadeer, a Uighur human
rights advocate who had been imprisoned in China and now lives in Washington, Xinhua reported. As
with the Dalai Lama, the Tibetan spiritual leader, Chinese officials often blame Ms.~Kadeer for
ethnic unrest; she denies the charges.

The clashes on Sunday began when the police confronted a protest march held by Uighurs to demand a
full government investigation of a brawl between Uighur and Han workers that erupted in Guangdong
Province overnight on June 25 and June 26. The brawl took place in a toy factory and left 2 Uighurs
dead and 118 people injured. The police later arrested a bitter ex-employee of the factory who had
ignited the fight by starting a rumor that 6 Uighur men had raped 2 Han women at the work site,
Xinhua reported.

There was also a rumor circulating on Sunday in Urumqi that a Han man had killed a Uighur in the
city earlier in the day, said Adam Grode, an English teacher living in the neighborhood where the
rioting took place.

``This is just crazy,'' Mr.~Grode said by telephone Sunday night. ``There was a lot of tear gas in
the streets, and I almost couldn't get back to my apartment. There's a huge police presence.''

Mr.~Grode said he saw a few Han civilians being harassed by Uighurs. Rumors of Uighurs attacking Han
Chinese spread quickly through parts of Urumqi, adding to the panic. A worker at the Texas
Restaurant, a few hundred yards from the site of the rioting, said her manager had urged the
restaurant workers to stay inside. Xinhua reported few details of the riot on Sunday night. It said
that ``an unknown number of people gathered Sunday afternoon'' in Urumqi, ``attacking passers-by and
setting fire to vehicles.''

Uighurs are the largest ethnic group in Xinjiang but are a minority in Urumqi, where Han Chinese
make up more than 70 percent of the population of two million or so. The Chinese government has
encouraged Han migration to the city and other parts of Xinjiang, fueling resentment among the
Uighurs. Urumqi is a deeply segregated city, with Han Chinese there rarely venturing into the Uighur
quarter.

The Uighur neighborhood is centered in a warren of narrow alleyways, food markets and a large
shopping area called the Grand Bazaar or the Erdaoqiao Market, where the rioting reached its peak on
Sunday.

Mr.~Grode, who lives in an apartment there, said he went outside when he first heard commotion
around 6 p.m. He saw hundreds of Uighurs in the streets; that quickly swelled to more than 1,000, he
said.

Police officers soon arrived. Around 7 p.m., protesters began hurling rocks and vegetables from the
market at the police, Mr.~Grode said. Traffic had ground to a halt. An hour later, as the riot
surged toward the center of the market, troops in green uniforms and full riot gear showed up, as
did armored vehicles. Chinese government officials often deploy the People's Armed Police, a
paramilitary force, to quell riots.

By midnight, Mr.~Grode said, some of the armored vehicles had begun to leave, but bursts of gunfire
could still be heard.

\section{Growing Presence in the Courtroom: Cellphone Data as Witness}

Mikhail Mallayev, who was convicted in March of murdering an orthodontist whose wife wanted him
killed during a bitter custody battle, stayed off his cellphone the morning of the shooting in
Queens. But afterward, he chatted away, unaware that his phone was acting like a tracking device and
would disprove his alibi -- that he was not in New York the day of the killing.

Darryl Littlejohn, a nightclub bouncer, made call after call on his cellphone as he drove from his
home in Queens to a desolate Brooklyn street to dump the body of Imette St. Guillen, the graduate
student he was convicted this month of murdering.

The pivotal role that cellphone records played in these two prominent New York murder trials this
year highlights the surge in law enforcement's use of increasingly sophisticated cellular tracking
techniques to keep tabs on suspects before they are arrested and build criminal cases against them
by mapping their past movements.

But cellphone tracking is raising concerns about civil liberties in a debate that pits public safety
against privacy rights. Existing laws do not provide clear or uniform guidelines: Federal wiretap
laws, outpaced by technological advances, do not explicitly cover the use of cellphone data to
pinpoint a person's location, and local court rulings vary widely across the country.

In one case that unsettled cellphone companies, a sheriff in Alabama told a carrier he needed to
track a cellphone in an emergency involving a child -- she turned out to be his teenage daughter,
who was late returning from a date.

For more than a decade, investigators have been able to match an antenna tower with a cellphone
signal to track a phone's location to within a radius of about 200 yards in urban areas and up to 20
miles in rural areas. Now many more cellphones are equipped with global-positioning technology that
makes it possible to pinpoint a user's position with much greater precision, down to a few dozen
yards.

To determine where a suspect's phone was in the past -- as in the Mallayev and Littlejohn cases --
investigators use company records that show a phone's approximate location at the beginning and end
of a call.

To track suspects in real time, law enforcement officials must ask a phone company to ``ping,'' or
send a signal to, a phone; for the effort to succeed, the phone must be turned on, though it does
not have to be in use. The police can then use a vehicle with signal-tracking equipment to narrow
down the location.

The frequency and ease with which law enforcement agencies access cellphone data to track people is
difficult to assess. Civil liberties groups recently obtained data from the Justice Department
through a lawsuit showing that in some jurisdictions, including New Jersey and Florida, courts often
allow federal prosecutors to track the location of cellphone users in real time without search
warrants.

Investigators seeking warrants must provide a judge with probable cause that a crime has been
committed. But investigators often obtain cell-tracking records under lower standards of judicial
review -- through subpoenas, which are granted routinely, or through an intermediate type of court
order based on an argument that the information requested would be relevant to an investigation.

In what would be the highest-level court decision on the issue so far, a federal appeals court in
Pennsylvania is expected to rule this summer on whether search warrants are required for the most
basic cellphone tracking data -- the electronic footprints that cellphone users leave behind in
company records, often without realizing it.

In March, Google announced that it would require search warrants before releasing GPS data that
pinpoints the movements of customers who use its mapping applications -- like Latitude, which lets
people see where their friends are -- on their phones.

But phone and Internet companies want Congress to clarify the laws so that they are clear about
their legal responsibilities.

Civil libertarians do not oppose using cellphone surveillance to solve crimes or save people in
emergencies, but they worry that the legal gray area is enabling it to happen without much scrutiny
or discussion.

``The cost of carrying a cellphone should not include the loss of one's personal privacy,'' said
Catherine Crump, a lawyer for the American Civil Liberties Union, which filed a lawsuit along with
the Electronic Frontier Foundation after the Justice Department did not respond to a Freedom of
Information request for data. Federal and local law enforcement officials argue that people who obey
the law have nothing to fear from cellphone tracking.

``Law enforcement has a responsibility to keep pace with the latest advances in technology in order
to improve its efficiency in combating crime,'' said Richard A.~Brown, the Queens district attorney,
whose office successfully prosecuted Mr.~Mallayev, adding that criminals are ``unknowingly
Twittering with law enforcement'' whenever they use their cellphones.

The data obtained by the civil liberties union provides a rare glimpse into crime-fighting
techniques that law enforcement agencies are reluctant to talk about. Since Sept.~12, 2001, federal
prosecutors in New Jersey have gained access to cellphone tracking information without warrants in
98 investigations resulting in 83 prosecutions.

Investigators have used cellphone tracking in a variety of ways: in Michigan to trace a fugitive
suspected of killing his wife to a wilderness park, and in Florida to seek a suspected serial killer
eventually killed in a police shootout.

Joseph A.~Pollini, who oversaw kidnapping investigations as lieutenant commander of the New York
Police Department's major case squad, said he used it routinely -- in one dramatic case tracing a
ransom call to a house in Flushing, Queens, where the victim was rescued and the kidnapper arrested.

Cellphone records quickly sharpened investigators' focus on Mr.~Littlejohn in the death of Ms.~St.
Guillen, whom he met in the bar where he worked as a bouncer. Records from Feb.~25, 2006, showed
that he called friends several times between 7 and 8 p.m. as his cellphone signal bounced off
antenna towers near his home in South Jamaica, Queens, and off towers along the Belt Parkway that
straddled the little-trafficked stretch of Fountain Avenue where an anonymous caller reported
finding Ms.~St. Guillen's body.

Records showed Mr.~Mallayev's cellphone moving south from New York to his home in Atlanta in the
hours after the killing of Daniel Malakov, the Queens orthodontist whose wife was convicted of
hiring Mr.~Mallayev to kill him. Mr.~Mallayev's son's cellphone was tracked sending signals half an
hour before the shooting to the tower closest to where Mr.~Malakov was killed.

The five largest wireless carriers receive hundreds of requests a month from law enforcement just
for real-time tracking, said Albert Gidari Jr., a Seattle lawyer who handles legal issues involving
cellphone tracking for several companies.

Civil libertarians say users whose phones have GPS-based services are unwittingly creating records
that could give the government easy access to their movements.

In the case being weighed in Pennsylvania, the Federal District Court in Pittsburgh ruled last year
that a search warrant was required even for historical phone location records, which the government
had requested to track a suspect in a drug case.

The decision upheld a magistrate judge's ruling that people have a reasonable expectation of privacy
regarding their physical location. Most Americans, the ruling said, do not know that their
cellphones create a record of their movements and ``would be appalled'' to learn that the government
can access it without showing probable cause. The Justice Department has appealed the case to the
United States Court of Appeals for the Third Circuit.

The civil liberties union and the Electronic Frontier Foundation are supporting the lower court's
decision, and say laws are needed, as the foundation puts it, to ``keep Big Brother out of your
pocket.''

\section{New Light in Old Bulbs}

SANTA ROSA, Calif. -- When Congress passed a new energy law two years ago, obituaries were written
for the incandescent light bulb. The law set tough efficiency standards, due to take effect in 2012,
that no traditional incandescent bulb on the market could meet, and a century-old technology that
helped create the modern world seemed to be doomed.

But as it turns out, the obituaries were premature.

Researchers across the country have been racing to breathe new life into Thomas Edison's light bulb,
a pursuit that accelerated with the new legislation. Amid that footrace, one company is already
marketing limited quantities of incandescent bulbs that meet the 2012 standard, and researchers are
promising a wave of innovative products in the next few years.

Indeed, the incandescent bulb is turning into a case study of the way government mandates can spur
innovation.

``There's a massive misperception that incandescents are going away quickly,'' said Chris Calwell, a
researcher with Ecos Consulting who studies the bulb market. ``There have been more incandescent
innovations in the last three years than in the last two decades.''

The first bulbs to emerge from this push, Philips Lighting's Halogena Energy Savers, are expensive
compared with older incandescents. They sell for \$5 apiece and more, compared with as little as 25
cents for standard bulbs.

But they are also 30 percent more efficient than older bulbs. Philips says that a 70-watt Halogena
Energy Saver gives off the same amount of light as a traditional 100-watt bulb and lasts about three
times as long, eventually paying for itself.

The line, for now sold exclusively at Home Depot and on Amazon.com, is not as efficient as compact
fluorescent light bulbs, which can use 75 percent less energy than old-style bulbs. But the Energy
Saver line is finding favor with consumers who dislike the light from fluorescent bulbs or are
bothered by such factors as their slow start-up time and mercury content.

``We're experiencing double-digit growth and we're continuing to expand our assortment,'' said Jorge
Fernandez, the executive who decides what bulbs to stock at Home Depot. ``Most of the people that
buy that bulb have either bought a C.F.L. and didn't like it, or have identified an area that
C.F.L.'s don't work in.''

For lighting researchers involved in trying to save the incandescent bulb, the goal is to come up
with one that matches the energy savings of fluorescent bulbs while keeping the qualities that many
consumers seem to like in incandescents, like the color of the light and the ease of using them with
dimmers.

``Due to the 2007 federal energy bill that phases out inefficient incandescent light bulbs beginning
in 2012, we are finally seeing a race'' to develop more efficient ones, said Noah Horowitz, senior
scientist with the Natural Resources Defense Council.

Some of the leading work is under way at a company called Deposition Sciences here in Santa Rosa.
Its technology is a key component of the new Philips bulb line.

Normally, only a small portion of the energy used by an incandescent bulb is converted into light,
while the rest is emitted as heat. Deposition Sciences applies special reflective coatings to
gas-filled capsules that surround the bulb's filament. The coatings act as a sort of heat mirror
that bounces heat back to the filament, where it is transformed to light.

While the first commercial product achieves only a 30 percent efficiency gain, the company says it
has achieved 50 percent in the laboratory. No lighting manufacturer has agreed yet to bring the
latest technology to market, but Deposition Sciences hopes to persuade one.

``We built a better mouse trap,'' said Bob Gray, coating program manager at Deposition Sciences.
``Now, we're trying to get people to beat a path to our door.''

With the new efficiency standards, experts predict more companies will develop specialized
reflective coatings for incandescents. The big three lighting companies -- General Electric, Osram
Sylvania and Philips -- are all working on the technology, as is Auer Lighting of Germany and
Toshiba of Japan.

And a wave of innovation appears to be coming. David Cunningham, an inventor in Los Angeles with a
track record of putting lighting innovations on the market, has used more than \$5 million of his
own money to develop a reflective coating and fixture design that he believes could make
incandescents 100 percent more efficient.

``There's enormous interest,'' Mr.~Cunningham said. ``All the major lighting companies want an
exclusive as soon as we demonstrate feasibility.''

Both Mr.~Cunningham and Deposition Sciences have been looking into the work of Chunlei Guo, an
associate professor of optics at Rochester University, who announced in May that he had used lasers
to pit the surface of a tungsten filament. ``Our measurements show that the treated filament becomes
twice as bright with the same power consumption,'' Mr.~Guo said.

And a physics professor at Rensselaer Polytechnic Institute, Shawn-Yu Lin, is also seeing improved
incandescent performance by using a high-tech, iridium-coated filament that recycles wasted heat.
``The technology can get up to six to seven times more efficient,'' Mr.~Lin said.

Despite a decade of campaigns by the government and utilities to persuade people to switch to
energy-saving compact fluorescents, incandescent bulbs still occupy an estimated 90 percent of
household sockets in the United States. Aside from the aesthetic and practical objections to
fluorescents, old-style incandescents have the advantage of being remarkably cheap.

But the cheapest such bulbs are likely to disappear from store shelves between 2012 and 2014, driven
off the market by the government's new standard. Compact fluorescents, which can cost as little as
\$1 apiece, may become the bargain option, with consumers having to spend two or three times as much
to get the latest energy-efficient incandescents.

A third technology, bulbs using light-emitting diodes, promises remarkable gains in efficiency but
is still expensive. Prices can exceed \$100 for a single LED bulb, and results from a government
testing program indicate such bulbs still have performance problems.

That suggests that LEDs -- though widely used in specialized applications like electronic products
and, increasingly, street lights -- may not displace incumbent technologies in the home any time
soon.

Given how costly the new bulbs are, big lighting companies are moving gradually. Osram will
introduce a new line of incandescents in September that are 25 percent more efficient. The bulbs
will feature a redesigned capsule with higher-quality gas inside and will sell for a starting price
of about \$3. That is less than the Philips product already on the market, but they will have
shorter life spans. G.E. also plans to introduce a line of household incandescents that will comply
with the new standards.

Mr.~Calwell predicts ``a lot more flavors'' of incandescent bulbs coming out in the future. ``It's
hard to be an industry leader in the crowded C.F.L field,'' he said. ``But a company could truly
differentiate itself with a better incandescent.''

\section{As Labor Secretary, Finding Influence in Her Past}

Hilda L.~Solis often recalls some advice her high school guidance counselor gave her mother: ``Your
daughter is not college material. Maybe she should follow the career of her older sister and become
a secretary.''

Telling that story recently at the Hunter College commencement in Manhattan, Ms.~Solis roared into
the microphone that she, the daughter of immigrants, did become a secretary -- the nation's labor
secretary. The crowd thundered with applause.

After surmounting many obstacles, Ms.~Solis, a former four-term congresswoman, has become the first
Hispanic woman to serve as a cabinet member. The third of seven children, she grew up in a modest
home near a giant landfill just east of Los Angeles. Her mother worked in a toy factory, her father
in a battery recycling plant, where he contracted lead poisoning.

Like President Obama and Judge Sonia Sotomayor, whose confirmation to the Supreme Court would make
her the first Hispanic justice, Ms.~Solis, 51, pulled herself up through her own drive and
intelligence. Now, she is hoping to use those qualities to reinvigorate the Labor Department, which
she said became a neglected, atrophied, often pro-business backwater under President George W.~Bush.

She has promised a vigorous campaign to combat workplace violations, after government auditors found
that the Bush Labor Department sometimes did not follow up on complaints of minimum wage, overtime
and child labor violations. She said she planned to hire 250 more investigators and conduct a
nationwide outreach program so that workers knew their rights and employers knew their obligations.

``There are so many people I knew when I was growing up who were not even paid the minimum wage,''
Ms.~Solis said. ``People wouldn't know where to go to lodge a complaint. And if you didn't speak
good English, forget it.''

Her goal of restoring morale and efficiency to the 17,000-employee department would not be easy, she
said, partly because so many dedicated people had quit in recent years. When she spoke at a church
in Miami in March, she pledged to increase workplace enforcement and used a much-applauded line:
``There's a new sheriff in town.'' The next day, John Sweeney, the president of the A.F.L.-C.I.O.,
gave her a plastic badge that read ``sheriff.''

But such enthusiastic support is far from universal, particularly among Republicans.

``She is anti-business and extremely pro-labor,'' said Ray Haynes, a Republican former state senator
in California who often clashed with Ms.~Solis when she was a lawmaker in Sacramento. ``From my
experience with Hilda, she has always been doctrinaire. Hilda does not hear the other side of the
story.''

Ms.~Solis said she planned to make the Labor Department a major player in fixing the pension system
and creating green jobs. And with unemployment climbing above 14.5 million, she has vowed to
strengthen job-training programs.

When Mr.~Obama offered her a cabinet post, she recalled, ``he said if I wanted to work for him,
you're going to be the voice for working families and organized labor.''

For her that seems to run in the family. At the dinner table, she said, her father, an immigrant
from Mexico, used to hold forth about factory conditions and how he, a Teamster shop steward, pushed
to improve them.

``My father spoke to management and fought for the workers' health and safety,'' she said. ``That is
something I bring to the Department of Labor. Those values my father shared with me.''

Her mother, an immigrant from Nicaragua who often stood 10 hours a day at the toy factory, was so
outspoken about working conditions that she would have been fired if the union had not protected
her, Ms.~Solis said.

Following her parents' example, Ms.~Solis has become a fighter -- for immigrants, workers,
minorities and women. ``I've always been about seeking social justice and combating discrimination
and racism,'' she said. ``I always wanted to stand up and fight for the underdog.''

Labor leaders said they hoped Ms.~Solis would push to enact their No.~1 objective, legislation that
would make it easier to unionize. At the time she was nominated to be labor secretary, Ms.~Solis was
the only member of Congress serving on the board of a pro-union group, American Rights at Work, a
move that some supporters said caused Republicans to delay her confirmation.

With the Obama administration focused on fighting the recession and overhauling the health care
system, she, like the president, has at least for now left it to Democratic senators to try to
muster the 60 votes needed to overcome a filibuster on the unionization bill.

Ms.~Solis majored in political science at California State Polytechnic University, Pomona, and then
obtained a master's degree in public administration at the University of Southern California. While
in graduate school, she interned at the White House under Jimmy Carter and soon landed a job with
the Office of Management and Budget.

Returning to California, she was hired to run a state program that helped poor students in the San
Gabriel Valley go to college. Friends urged her to run for office, and at age 28, she was elected to
the board of Rio Hondo Community College.

In 1992, she won a seat in the California State Assembly and two years later she was elected to the
State Senate, becoming California's first Latina state senator.

Remembering her malodorous hometown, La Puente, where there is a federal Superfund site and eight
landfills, she sponsored a landmark environmental bill, which required the State Environmental
Protection Agency to adopt regulations ensuring fair treatment of people of all races and incomes
with respect to environmental laws.

Impressed by the legislation, the John F.~Kennedy Library Foundation awarded her the Profile in
Courage Award in 2000; she was the first woman to win it.

``This legislation was groundbreaking,'' Caroline Kennedy said in a telephone interview. ``The bill
didn't succeed the first time, but she came back and worked with the business community,
compromised, and ultimately succeeded.''

That same year, Ms.~Solis challenged a nine-term congressman, Matthew Martinez, in the Democratic
primary, angering many Hispanic leaders. Even though Mr.~Martinez had a generally liberal voting
record, she accused him of being out of touch with his constituency.

Unhappy about Mr.~Martinez's free-trade voting record, the Los Angeles County Federation of Labor
endorsed her, with its president calling her a ``warrior for working families.'' She won the primary
62 percent to 29 percent and ran unopposed that November.

One of her mentors, Dolores Huerta, a founder of the United Farm Workers, said she was mystified by
one thing about Ms.~Solis. ``I often went to Washington, and I always wondered why Hilda never said,
'Stay at my place,' '' Ms.~Huerta said. ``One night we had dinner, and she had to stop at her
apartment to pick up something. Her apartment was tiny. There was barely enough room for a bed. I
thought, 'This is how a nun lives.' ''

Ms.~Solis is married; her husband, Sam Sayyad, runs an auto repair business outside Los Angeles.

At the Hunter commencement, Ms.~Solis concluded by returning to her own story.

``People always say that women, people of color, Latinas, they're not ready to go to college,
they're not ready to be in those big positions,'' she said. ``There are probably a dozen of you in
this hall who are future Sonia Sotomayors, and there are probably two dozen future Hilda Solises.
You have to have the ganas -- the desire -- to do it.''

\section{President's Vacation: Classified, but Not Much of a Secret}

WASHINGTON -- He does not own a sprawling Texas ranch, a family compound on the coast of Maine or a
retreat tucked away in the Santa Ynez Mountains of California.

So President Obama, in this regard at least, is much like Bill Clinton: He has no vacation home to
call his own.

With the summer getaway season in full swing, the White House is busy arranging the president's
first vacation since taking office. The destination -- officially, at least -- is classified. Yet it
is hardly a secret to the people on Martha's Vineyard, several business owners and others said,
where reservations have been made and preparations are under way for the Obama family's August
arrival.

Still, questions abound on the Massachusetts island.

Will they settle in Oak Bluffs, Chilmark or another town? Will they rent the home of the filmmaker
Spike Lee? Will they be there at the same time as the Clintons?

``The rumor mill is rampant because people want to get information and want to flesh out their
plans,'' said Charles J.~Ogletree, a law professor at Harvard, a summer Vineyard resident and a
longtime friend of the Obamas. ``I hope that people will respect their privacy and need for rest.
The best thing we can do is let him recharge his batteries and prepare for the tasks ahead.''

The White House is mum on the dates, duration and precise destination of the trip.

First, advisers said, the administration does not want to distract from the work of Congress, which
will spend at least some of the summer grinding away on health care legislation and the Supreme
Court confirmation hearings for Judge Sonia Sotomayor.

Second, as the unemployment rate inches ever closer to 10 percent, an oceanfront respite filled with
tee times and dinner dates could be seen as out-of-touch with the economic hardship facing many
Americans.

And third, the president is really never off duty, so wherever he chooses to vacation -- and for how
long -- sends at least some type of a message and is sure to draw scrutiny.

Last August, when Mr.~Obama flew to his native Hawaii, some Washington commentators criticized his
vacation choice. Mr.~Obama, who also spent time visiting his dying grandmother in Honolulu, told
aides he was particularly agitated by a remark from Cokie Roberts of ABC News, who said, ``It has
the look of him going off to some sort of foreign, exotic place.'' She added, ``He should be in
Myrtle Beach.''

This year, aides said, a handful of destinations were discussed. (Myrtle Beach was not among them.)
But since the Obamas have previously visited Martha's Vineyard, and the Secret Service is accustomed
to protecting presidents there, it became a top choice.

The island, which is 23 miles long and 9 miles wide, has six towns with distinctive traits. The
Obamas are likely to stay in or near Oak Bluffs, where upper-class black families have been
vacationing for decades.

``Vineyarders allow people to enjoy their time here without bothering them,'' said Nancy Gardella,
executive director of the Martha's Vineyard Chamber of Commerce. ``They pride themselves on kind of
ignoring celebrities and let them relax without being hassled or followed.''

That does not mean, of course, that Mr.~Obama will not be hassled from afar.

During his Hawaiian trip last summer, the Republican National Committee carefully tracked his
itinerary, sending out pictures and news accounts of his whereabouts on Oahu. This year, the themes
of criticism could well come from the words of previous Republican presidents.

``Most Americans don't sit in Martha's Vineyard, swilling white wine,'' George W.~Bush once said
when asked why he chose to relax in the broiling Texas sun as opposed to the cool Atlantic breezes.
And Lyndon B.~Johnson disparagingly referred to the Vineyard as the ``female island,'' upon
discovering that his aides chose Martha's Vineyard for their vacation instead of joining him at his
ranch in Texas.

The full-time residents of the island pay little attention to the jokes and are eager to welcome the
Obamas to a place that has also felt the sting of the depressed economy.

``Having the Obama family pick us is great for business,'' said Susan Goldstein, the owner of the
Mansion House hotel, who has lived on the Vineyard for 40 years. ``It's like our own islandwide
recovery package.''

In 1993, Ms.~Goldstein hung a welcome banner when the Clintons arrived for a 10-day vacation. She
said that she planned to do the same for the Obama family, but that she suspected the presidential
trips would be entirely different. On one of Mr.~Clinton's summer visits, he attended nine cocktail
or dinner parties in 13 nights, captivating the island social scene.

Mr.~Obama and his wife, Michelle, are planning a far more low-key visit, several aides and
associates say. A few Chicago friends are expected to join the vacation, as they have trips to
Hawaii. And golf outings on any of the island's five courses -- two public, three private -- are
expected to be with friends, rather than the celebrity partners Mr.~Clinton played with.

``It won't be a constant series of public events, but a genuine opportunity to spend quality time
with their daughters and some important time with a group of friends,'' Mr.~Ogletree said. ``I hope
and pray this is not a working vacation, where they don't have any time for themselves, and go back
to Washington more exhausted than they came. That would be a tragic and unfortunate intrusion on
their time.''

While Mrs.~Obama and the couple's two daughters, Malia and Sasha, are planning to stay longer, aides
said Mr.~Obama was likely to take only a week of vacation. The Obama administration is intent on
avoiding any parallels to Mr.~Bush, who spent 27 days at his Texas ranch during the first August
vacation of his presidency, which came only days before the terrorist attacks of Sept.~11, 2001.

But unlike Mr.~Bush, who established a Western White House in Crawford and invited members of his
cabinet and his senior aides to join him, Mr.~Obama intends to be away from most of his team. Rahm
Emanuel, the president's chief of staff, is heading to Montana. David Axelrod, a senior adviser, is
going to his lakefront house in Michigan. Valerie Jarrett, another senior adviser, will be on the
Vineyard but at a house she has rented for years.

With the exception of a few weekend trips to Camp David, the presidential retreat in the Catoctin
Mountains of Maryland, Mr.~Obama has taken no vacation in his first five months in office. Last
year, he waxed aloud about the importance of getting away.

``The most important thing you need to do is to have big chunks of time during the day when all
you're doing is thinking,'' Mr.~Obama said. If not, he said, ``you start making mistakes or you lose
the big picture.''

\end{document}




